\chapter{量子修正与规范结构}\label{chap:quantum-radiative-structure}

树级振幅只使用有限条内部传播线；闭合圈则要求积分所有允许的内部动量。高动量部分可能发散，中间态达到质量壳又会产生虚部。前者关乎参数怎样由实验定义，后者关乎散射概率怎样守恒，两种问题必须从同一个振幅中分别辨认。

幺正性把虚部与真实中间态联系起来，局域重整化把短距离发散吸收到允许的参数与算符中。规范场的冗余还限制哪些传播和顶角修正能够共同出现；若真空具有简并结构，其附近的涨落则决定新的低能自由度。具体规范理论将在这些条件下计算，而不是先给圈图指定一个孤立的有限答案。

\section{量子修正与规范结构}

树级振幅只包含有限条内部传播线；闭合圈出现以后，内部四动量必须在所有可能取值上积分，于是首先要回答两个具体问题：这些积分的高动量部分怎样定义，以及中间态真正达到质量壳时为什么会给振幅带来虚部。规范场还包含并非独立物理自由度的冗余分量；具有简并最低能态的理论则在选定真空后出现新的低能自由度。

$S$ 矩阵的幺正性首先把振幅虚部与真实中间态联系起来，这使圈积分的解析结构具有直接的概率解释。圈积分的紫外部分又迫使裸参数改写成由实验条件定义的有限参数。规范场还多出一个结构性问题：规范冗余使二次核不可逆，必须先选定规范截面才能定义传播子。若势能最低点不是孤立点，真空附近的二次涨落还会沿简并方向出现新的低能结构。

\subsection{幺正性与割线}
散射概率不能在时间演化中凭空增加或减少，因此 $S$ 矩阵必须满足幺正性。把这一条件展开到给定初态和末态之间，会发现振幅的虚部并不是独立信息：它由所有允许的真实在壳中间态共同决定。等到这个算符关系写成具体圈图积分时，才得到通常所称的 Cutkosky 割线规则。
% 幺正性、支割线与 Cutkosky 关系
圈图中的内部传播子通常不在质量壳上，因此它们只是积分变量，并不对应可直接探测的中间粒子。随着外部总能量升高，某些内部线可能第一次同时满足正能质量壳条件；从这一阈值开始，积分中出现与真实中间态相同的相空间。Feynman 分母中的 $+\ii0^+$ 规定决定这个新贡献以怎样的虚部进入振幅。把“哪些内部线同时在壳”画成切开图以后，才得到通常所说的 Cutkosky 规则；所谓极点夹逼（pinch）只是这一阈值在复积分平面上的解析表现。


二体阈值的解析结构已经包含在两个标量传播子的共同分母中，与具体顶角分子无关。取完成 $k^0$ 轮廓积分后的标量泡图核
\begin{equation}
 \mathcal B(P^0)
 \equiv
 \int\frac{\dd^3 \mathbf k}{(2\pi\hbar)^3}
 \frac{c^2}{4E_1E_2}
 \left[
 \frac{1}{P^0c-E_1-E_2+\ii0^+}
 -\frac{1}{P^0c+E_1+E_2-\ii0^+}
 \right],
 \label{eq:qft-bubble-after-k0-dp8}
\end{equation}
其中
\begin{equation*}
 E_1=\sqrt{m_1^2c^4+c^2\mathbf k^2},
 \qquad
 E_2=\sqrt{m_2^2c^4+c^2(\mathbf P-\mathbf k)^2}.
\end{equation*}
对正的 $P^0$，只有第一项的分母能够在实积分域取零，因此其虚部为
\begin{align}
 \operatorname{Im}\mathcal B(P^0)
 &=-\pi
 \int\frac{\dd^3 \mathbf k}{(2\pi\hbar)^3}
 \frac{c^2}{4E_1E_2}
 \delta(P^0c-E_1-E_2).
 \label{eq:bubble-imag-phase-space-dp8}
\end{align}
这个 $\delta$ 函数正是“两个内部粒子同时成为真实正能粒子”的能量守恒条件。对等质量情形 $m_1=m_2=m$，质心系中的在壳动量和速度因子为
\begin{align}
 k_*^2
 &=\frac{s-4m^2c^2}{4},\\
 \beta_*(s)
 &\equiv\frac{ck_*}{E_{\rm cm}/2}
 =\sqrt{1-\frac{4m^2c^2}{s}}.
 \label{eq:two-body-threshold-beta-dp8}
\end{align}
因此二粒子连续谱只能在 $s\ge4m^2c^2$ 后打开，并在阈值附近按 $\beta_*(s)$ 从零增长。

传播子割线的局部规则来自同一个分布恒等式。先写单个 Feynman 分母跨过实轴的差：
\begin{equation}
 \operatorname{Disc}\frac1{q^2-m^2c^2+\ii0^+}
 =-2\pi\ii\,\delta(q^2-m^2c^2).
 \label{eq:feynman-denominator-discontinuity-dp68}
\end{equation}
物理中间态还必须选取正能支撑，于是图中被切开的标量传播子采用
\begin{equation}
 \boxed{
 \frac1{q^2-m^2c^2+\ii0^+}
 \quad\longrightarrow\quad
 -2\pi\ii\,\delta_+(q^2-m^2c^2)}
 \label{eq:cutkosky-propagator-rule-dp8}
\end{equation}
其中 $\delta_+(q^2-m^2c^2)\equiv\Theta(q^0)\delta(q^2-m^2c^2)$。若采用带显式 $\ii$ 的传播子定义，相应整体 $\ii$ 因子也必须按同一约定保留。

另一方面，前文散射幺正性正文\eqnrefs{eq:optical-unitarity-r68} 可以按耦合阶数逐阶展开。最低非平凡阶给出
\begin{equation}
 \boxed{
 2\operatorname{Im}\langle i|\hat T^{(2)}|i\rangle
 =\sum_X
 \langle i|\hat T^{(1)\dagger}|X\rangle
 \langle X|\hat T^{(1)}|i\rangle,}
 \label{eq:perturbative-optical-theorem-dp68}
\end{equation}

\begin{proposition}[幺正性的逐阶恒等式]
设 $S=I+\ii T$ 幺正，并按同一个实微扰参数展开；逐阶相等按形式幂级数理解。则有以下结论（原式 \eqnrefs{eq:perturbative-optical-theorem-dp68}）：
\begin{equation*}
 \boxed{
 2\operatorname{Im}\langle i|\hat T^{(2)}|i\rangle
 =\sum_X
 \langle i|\hat T^{(1)\dagger}|X\rangle
 \langle X|\hat T^{(1)}|i\rangle,}
\end{equation*}
\end{proposition}
\begin{proof}
正文\eqnrefs{eq:optical-unitarity-r68} 是精确散射算符的幺正性关系。把跃迁算符按耦合阶数展开，并只比较最低出现虚部的二阶项，就能把圈振幅的虚部写成一阶真实跃迁振幅的模平方之和，从而得到正文\eqnrefs{eq:perturbative-optical-theorem-dp68}。

写
\begin{equation*}
 \hat T_S=\hat T^{(1)}+\hat T^{(2)}+\hat T^{(3)}+\cdots .
\end{equation*}
代入
$-\ii(\hat T_S-\hat T_S^\dagger)=\hat T_S^\dagger\hat T_S$，二阶部分为
\begin{equation*}
 -\ii\bigl(\hat T^{(2)}-\hat T^{(2)\dagger}\bigr)
 =\hat T^{(1)\dagger}\hat T^{(1)}.
\end{equation*}
在同一初态 $|i\rangle$ 两侧取矩阵元，并在右端插入在壳中间态完备关系 $\sum_X|X\rangle\langle X|=\mathbb I$，得到
\begin{align*}
 -\ii\left[
 \langle i|\hat T^{(2)}|i\rangle
 -\langle i|\hat T^{(2)}|i\rangle^*
 \right]
 &=\sum_X
 \langle i|\hat T^{(1)\dagger}|X\rangle
 \langle X|\hat T^{(1)}|i\rangle\\
 &=2\operatorname{Im}\langle i|\hat T^{(2)}|i\rangle,
\end{align*}
这就是正文\eqnrefs{eq:perturbative-optical-theorem-dp68}。图形语言中的割线规则只是把右端的在壳中间态完备和翻译成内部传播子的 $\delta_+$ 替换。
\end{proof}

其中 $|X\rangle$ 是完整的真实在壳中间态。\eqnrefs{eq:cutkosky-propagator-rule-dp8} 正是把这一算符关系翻译成圈积分的逐图计算规则。

\begin{figure}[htbp]
 \centering
 \includegraphics[width=0.78\textwidth]{two_body_cut_threshold.png}
 \caption{等质量二粒子割线的无量纲阈值因子 $\beta_*(s)$，由\eqnrefs{eq:two-body-threshold-beta-dp8} 直接计算。横轴 $x=\sqrt{s}/(2mc)$；$x=1$ 是二粒子阈值。该图只展示相空间开启因子，具体圈图的整体前因子和分子结构由相互作用决定。}
 \label{fig:two-body-cut-threshold-dp8}
\end{figure}

\begin{figure}[htbp]
 \centering
 \begin{tikzfig}[x=1cm,y=1cm]
 \coordinate (L) at (-2.2,0); \coordinate (R) at (2.2,0);
 \coordinate (a) at (-0.9,0); \coordinate (b) at (0.9,0);
 \draw[wfA] (L)--(a);
 \draw[wfA] (b)--(R);
 \draw[wfA] (a) .. controls (-0.25,1.25) and (0.25,1.25) .. (b);
 \draw[wfA] (a) .. controls (-0.25,-1.25) and (0.25,-1.25) .. (b);
 \fill (a) circle (1.7pt); \fill (b) circle (1.7pt);
 \draw[dashed,wfA] (0,-1.55)--(0,1.55);
 \node at (-1.65,.3) {$P$};
 \node[anchor=south east,fill=white,inner sep=1pt] at (-.12,1.42) {$k$};
 \node[anchor=north west,fill=white,inner sep=1pt] at (.12,-1.42) {$P-k$};
 \node[align=center,right] at (2.7,0) {$k^2=m_1^2c^2$\\$(P-k)^2=m_2^2c^2$};
\end{tikzfig}

 \caption{二体泡图的幺正性割线。虚线表示把两条内部传播子同时置于正能质量壳；割线积分使用的测度就是相应真实二体 Lorentz 不变相空间。}
 \label{fig:qft-bubble-cut-dp8}
\end{figure}





\subsection{紫外发散、尺度依赖与短距离展开}
圈积分会探测任意大的内部动量，因此中间表达式可能依赖一个人为引入的紫外调节尺度。物理预测不能依赖这个人为选择：需要先把发散写成局域参数的重新定义，再研究当参考尺度改变时这些有限参数怎样随之变化。这种有限参数随重整化尺度的演化就是重整化群。

当计算推广到多圈或局域复合算符时，同一个“短距离效应应由局域结构吸收”的原则仍然成立，只是需要更系统的递归组织。等真正比较两个相近时空点上的算符乘积时，再引入算符乘积展开；热核则只在需要计算一般二次算符的一圈紫外项时出现。

分母合并、欧氏径向积分和维数极点的展开已分别由\eqnrefs{eq:feynman-two-general-r68,eq:feynman-three-denominator-r14,eq:euclidean-master-integral-dp21,eq:dimreg-universal-expansion-dp68}给出。现在要把其中的紫外极点与可由实验确定的有限参数联系起来。
\phantomsection\label{sec:qft-renormalization-rg-dp10}

前面的微扰规则已经把散射振幅写成顶角、传播子和内部动量积分的组合。新的困难出现在圈积分的高动量端：某些积分若一直延伸到任意大的内部动量，会发散，因而“直接代入裸质量和裸耦合再积分”并不是一个完整的计算定义。解决它需要先把发散积分改写成带有临时控制参数的有限表达式，再用实验可定义的质量、耦合和场归一化重新参数化作用量，使最终可观测量不依赖这个临时参数。前一步称为正则化，后一步称为重整化。

重整化以后，同一个物理过程仍可在不同外部能标下被探测，因此有限参数必须随分辨尺度变化而保持可观测量不变；描述这种尺度依赖的微分方程就是重整化群。当实验能量远低于某些重自由度的质量尺度时，这些自由度不能成为真实外态，却仍通过虚过程留下低能效应；有效场论把这些高能效应编码成局域算符及其 Wilson 系数。

圈动量定义为顶角守恒约束后仍未固定的独立内部 Fourier 积分变量；一般内部传播线处于离壳状态。紫外发散说明“裸参数直接等于实验参数”的写法没有把短程贡献与测量条件分开。正则化引入临时数学参数；反项与重整化条件再把这种临时依赖吸收到有限个参数定义中。最终可观测量只依赖物理输入与外部运动学尺度，而不依赖正则化器本身。


共同放大所有独立圈物理四动量，\(\ell_r^\mu\mapsto\Lambda\ell_r^\mu\) 且 \(\Lambda\to\infty\)，即可诊断图的表面紫外次数。四维圈测度贡献四次幂，玻色传播子贡献负二次幂，Dirac 传播子在大动量下贡献负一次幂；若顶角分子还含 \(N_\partial\) 个圈动量因子，则连通图满足
\begin{equation}
 {
 \omega_{\rm UV}=4L-2I_B-I_F+N_\partial,}
 \label{eq:superficial-degree-general-dp10}
\end{equation}
其中 \(L\) 是独立圈数，\(I_B\) 与 \(I_F\) 分别为内部玻色子线与费米子线数。\(\omega_{\rm UV}<0\) 只说明所有圈动量共同放大时表面收敛，仍须检查发散子图；\(\omega_{\rm UV}\ge0\) 也只是最坏幂次上界，对称性还可以排除相应局域结构。

对四维 \(\varphi^4\) 理论，每个四次顶角提供四个线端。若连通图含 $V$ 个顶角、$I$ 条内部线和 $E$ 条外线，则线端计数与独立圈数满足
\begin{equation}
 {
 4V=2I+E,
 \qquad
 L=I-V+1.}
 \label{eq:phi4-topology-r68}
\end{equation}
把这两条纯拓扑关系代入\eqnrefs{eq:superficial-degree-general-dp10}，便得到
\begin{equation}
 {\omega_{\varphi^4}=4-E,}
 \label{eq:phi4-superficial-degree-dp10}
\end{equation}
所以除真空图外，只有二点与四点函数需要原理论已有的质量、场强和四次耦合局域结构作为反项。QED 的每个最小耦合顶角连接一个光子腿和两个费米子腿。相应的光子线端、费米子线端和独立圈数同时满足
\begin{equation}
 {
 V=2I_\gamma+E_\gamma,
 \qquad
 2V=2I_\psi+E_\psi,
 \qquad
 L=I_\gamma+I_\psi-V+1.}
 \label{eq:qed-topology-r68}
\end{equation}
把这些拓扑恒等式代入 $\omega=4L-2I_\gamma-I_\psi$，即可消去内部线数并得到
\begin{equation}
 {
 \omega_{\rm QED}=4-E_\gamma-\frac32E_\psi.}
 \label{eq:qed-superficial-degree-dp10}
\end{equation}
例如光子二点函数表面计数给二次发散，但 Ward 恒等式要求 \(q_\mu\Pi^{\mu\nu}=0\)，把允许张量限制到横向结构，实际只剩对数发散。幂次计数决定潜在紫外阶数，规范对称性再决定哪些局域项真正允许出现。

发散积分必须先被数学地定义。维数正则化取
\begin{equation*}
 d=4-2\epsilon,
 \qquad \epsilon>0,
\end{equation*}
并引入具有\emph{物理四动量量纲}的尺度 \(\mu_{\rm DR}\)。四维无量纲耦合的一般圈积分写成
\begin{equation}
 {
 \mu_{\rm DR}^{2\epsilon}
 \int\frac{\dd^d \ell}{(2\pi)^d}(\cdots),}
 \label{eq:dim-reg-measure-general-dp10}
\end{equation}
其中 $\ell^\mu$ 始终是物理四动量。若改用逆长度变量 $\kappa^\mu=\ell^\mu/\hbar$，则测度提供 $\hbar^{-d}$，两个二次传播子分母提供 $\hbar^4$，而维数正则化尺度 $\mu_\kappa^{2\epsilon}=(\mu_{\rm DR}/\hbar)^{2\epsilon}$ 再提供 $\hbar^{-2\epsilon}$。因为 $d=4-2\epsilon$，总幂次
\[
 -d+4-2\epsilon=0
\]
恰好为零，因此物理动量形式中的显式 $\hbar$ 由变量换元严格抵消。典型积分在 $\epsilon\to0$ 时产生局域极点与外部尺度对数；前者由反项吸收，后者保留为可测尺度依赖。

有外部四动量时，一圈分子通常还含 $\ell^\mu$、$\ell^\mu\ell^\nu$ 等张量。Lorentz 协变性保证它们只能展开在 $\eta^{\mu\nu}$ 与外部四动量张量基上，因此张量积分可以先约化为少数标量母积分；这一步称为 Passarino--Veltman 约化。沿用\eqnrefs{eq:reduced-loop-measure-si-r10} 的物理四动量测度，对任意具有动量量纲的质量参数 $M_i$ 定义
\begin{equation}
{
\begin{aligned}
 A_0(M^2)
 &\equiv \frac{16\pi^2}{\ii}\int_\ell^{(d)}
 \frac{1}{\ell^2-M^2+\ii0},\\
 B_0(p^2;M_1^2,M_2^2)
 &\equiv \frac{16\pi^2}{\ii}\int_\ell^{(d)}
 \frac{1}{(\ell^2-M_1^2+\ii0)[(\ell+p)^2-M_2^2+\ii0]},\\
 C_0(p_1^2,p_2^2,(p_1+p_2)^2;M_1^2,M_2^2,M_3^2)
 &\equiv \frac{16\pi^2}{\ii}\int_\ell^{(d)}
 \frac{1}{D_1D_2D_3},
\end{aligned}}
\label{eq:pv-scalar-masters-dp90}
\end{equation}
其中
\[
D_1=\ell^2-M_1^2+\ii0,\quad
D_2=(\ell+p_1)^2-M_2^2+\ii0,\quad
D_3=(\ell+p_1+p_2)^2-M_3^2+\ii0.
\]
按这一定义，$A_0$、$B_0$、$C_0$ 的动量量纲依次为 $2,0,-2$。例如秩二三点积分不需要引入新的非标量函数作为物理输入，因为必可写成
\begin{equation}
\frac{16\pi^2}{\ii}\int_\ell^{(d)}\frac{\ell^\mu\ell^\nu}{D_1D_2D_3}
=C_{00}\eta^{\mu\nu}+\sum_{i,j=1}^{2}C_{ij}p_i^\mu p_j^\nu .
\label{eq:pv-rank2-decomposition-dp90}
\end{equation}
分别用 $\eta_{\mu\nu}$、$p_{1\mu}p_{1\nu}$、$p_{1\mu}p_{2\nu}$、$p_{2\mu}p_{2\nu}$ 收缩即可得到 $C_{00},C_{ij}$ 的线性方程；再利用 $2\ell\!\cdot p_i$ 与相邻传播子分母之差，把右端化为\eqnrefs{eq:pv-scalar-masters-dp90}。因此 QED、QCD 与标准模型的一圈张量约化都归结为同一件事：先在 Lorentz 张量基上展开，再用传播子恒等式把张量积分精确消元为标量主积分。

重整化把同一裸作用量精确拆成“用有限参数写成的原型拉氏量 + 局域反项”。以实标量场为例，写成
\begin{equation*}
 \varphi_B=Z_\varphi^{1/2}\varphi_R,
 \qquad
 m_B^2=m_R^2+\delta m^2+\cdots,
 \qquad
 g_{4,B}=\mu_{\rm DR}^{2\epsilon}(g_4+\delta g_4),
\end{equation*}
即可得到
\begin{equation*}
 \Lag_B=\Lag_{\rm R}+\Lag_{\rm ct}.
\end{equation*}
反项系数由重整化条件以及裸理论对临时正则化参数的独立性确定。

\begin{note}[常用重整化方案]
壳上方案以传播子物理极点、留数和指定运动学点的电荷作为输入；最小减除只减去 \(1/\epsilon\) 极点；\(\overline{\rm MS}\) 方案同时减去 \(1/\epsilon-\gamma_E+\ln4\pi\)。有限阶计算会保留方案依赖，但同一可观测量在相同完整微扰阶上的方案差异属于更高阶截断误差。
\end{note}


\begin{exbox}[$\varphi^4$ 理论的三类重整化]
单靠“圈积分发散”还看不出应该重定义哪个参数。最清楚的办法是分别看二点函数与四点函数。

以四维 $\varphi^4$ 模型为对象，比较二点与四点圈图需要的反项，确定质量、场强和耦合的重整化条件。
\textbf{解。} 一圈二点 1PI 图只有一个四次顶角。两条外腿占用顶角的两个场，剩余两个场彼此收缩成蝌蚪圈，因此
\begin{equation}
 {
 -\ii\Sigma_{\rm tad}^{(1)}
 =\frac{-\ii g_4}{2}
 \mu_{\rm DR}^{2\epsilon}
 \int\frac{\dd^d \ell}{(2\pi)^d}
 \frac{\ii}{\ell^2-m_R^2c^2+\ii0^+}.}
 \label{eq:phi4-tadpole-selfenergy-dp52}
\end{equation}
$1/2$ 是同一顶角上两条内部腿互换的对称性因子。这个积分不含外部四动量 $p$，所以它只能平移传播子极点的位置，而不能改变极点附近对 $p^2$ 的斜率。于是四维 $\varphi^4$ 理论在一圈阶满足
\begin{equation}
 {
 \delta m^2=O(g_4),
 \qquad
 Z_\varphi=1+O(g_4^2).}
 \label{eq:phi4-one-loop-mass-vs-Z-dp52}
\end{equation}
也就是说，质量反项在一圈就需要，场强反项要到二圈才第一次被二点函数要求。

四点函数的情况不同。$g_4^2$ 阶出现 $s,t,u$ 三个泡图，每个都依赖外部不变量，并具有同一个局域对数紫外极点。它们不能由质量反项消去，而必须重定义四次耦合。于是“哪一种反项出现”并不是预先背下来的规则，而是由 1PI 函数的外腿数和外部动量结构决定：二点函数的常数部分固定质量，二点函数的 $p^2$ 斜率固定场强，四点函数在指定运动学点的值固定耦合。

把重整化后二点函数写成
\begin{equation}
 \widetilde G_R(p)
 =\frac{\ii}
 {p^2-m_R^2c^2-\Sigma_R(p^2)+\ii0^+},
 \label{eq:phi4-ren-propagator-dp52}
\end{equation}
壳上方案直接把“物理质量”和“单位极点留数”翻译成
\begin{equation}
 {
 \Sigma_R(m_R^2c^2)=0,
 \qquad
 \left.\frac{\dd\Sigma_R}{\dd p^2}\right|_{p^2=m_R^2c^2}=0.}
 \label{eq:phi4-onshell-ren-conditions-dp52}
\end{equation}
第一式固定极点位置，第二式固定极点斜率。$\overline{\rm MS}$ 不直接采用这两个物理条件，而只减去维数正则化极点；两种方案得到的中间参数不同，但把参数都换成同一组实验输入后，可观测量只在被截掉的更高微扰阶上不同。


\medskip

\eqnrefs{eq:phi4-tadpole-selfenergy-dp52} 是 $\varphi^4$ 二点 1PI 函数的一圈蝌蚪贡献。由于积分与外部四动量无关，它在这一阶只移动传播子的极点位置，而不产生随 $p^2$ 变化的紫外斜率项。顶角收缩计数给出 $1/2$ 对称因子；维数正规化把紫外极点显式分离，质量与场强反项分别控制逆传播子的常数项和 $p^2$ 斜率。把重整化逆传播子在物理极点附近展开并固定极点位置与留数，便得到\eqnrefs{eq:phi4-onshell-ren-conditions-dp52}。



从相互作用
\begin{equation*}
 \Lag_{\rm int}
 =-\frac{g_4}{4!\,\hbar c}\varphi^4
\end{equation*}
的一阶 Dyson 插入开始。为形成二点 1PI 函数，两个外部场必须分别与顶角上的两个场收缩；选择有序外腿的方式数为
\begin{equation*}
 4\times3=12.
\end{equation*}
剩余两个顶角场彼此收缩只有一种方式。作用量中的 $1/4!=1/24$ 因而留下
\begin{equation*}
 \frac{12}{24}=\frac12,
\end{equation*}
这就是\eqnrefs{eq:phi4-tadpole-selfenergy-dp52} 的对称性因子。内部线从顶角出发又回到同一顶角，没有任何外部四动量流经圈，因此积分只依赖 $m_R$ 与正则化尺度，而不依赖 $p$。

令
\begin{equation*}
 T(m_R^2)
 \equiv
 \mu_{\rm DR}^{2\epsilon}
 \int\frac{\dd^d \ell}{(2\pi)^d}
 \frac{\ii}{\ell^2-m_R^2c^2+\ii0^+}.
\end{equation*}
用 Schwinger 参数
\begin{equation*}
 \frac{\ii}{\ell^2-m_R^2c^2+\ii0^+}
 =\int_0^\infty\dd\tau\,
 \ee^{\ii\tau(\ell^2-m_R^2c^2+\ii0^+)},
\end{equation*}
把动量积分化成 $d$ 维高斯积分，再对 $\tau$ 积分，得到解析延拓形式
\begin{equation*}
 T(m_R^2)
 =\frac{\mu_{\rm DR}^{2\epsilon}}
 {(4\pi)^{d/2}}
 \Gamma\!\left(1-\frac d2\right)
 (m_R^2c^2)^{d/2-1}.
\end{equation*}
取 $d=4-2\epsilon$，并展开
\begin{align*}
 \Gamma(-1+\epsilon)
 &=-\frac1\epsilon+\gamma_E-1+O(\epsilon),\\
 (4\pi)^\epsilon
 \left(\frac{\mu_{\rm DR}^2}{m_R^2c^2}\right)^\epsilon
 &=1+\epsilon
 \left[
 \ln4\pi+\ln\frac{\mu_{\rm DR}^2}{m_R^2c^2}
 \right]+O(\epsilon^2),
\end{align*}
所以
\begin{equation*}
 T(m_R^2)
 =-\frac{m_R^2c^2}{16\pi^2}
 \left[
 \frac1{\bar\epsilon}
 +1-\ln\frac{m_R^2c^2}{\mu_{\rm DR}^2}
 \right]
 +O(\epsilon),
\end{equation*}
其中 $1/\bar\epsilon=1/\epsilon-\gamma_E+\ln4\pi$。这个结果的关键结构是“局域、与 $p$ 无关、正比于 $m_R^2c^2$”；其有限常数部分会随减除方案改变，而这三个结构不会改变。

把裸二次部分写成重整化场与反项：
\begin{equation*}
 \frac12 Z_\varphi(\partial\varphi_R)^2
 -\frac12 Z_\varphi m_B^2c^2\frac{\varphi_R^2}{\hbar^2}
 =\frac12(\partial\varphi_R)^2
 -\frac12m_R^2c^2\frac{\varphi_R^2}{\hbar^2}
 +\Lag_{\rm ct}^{(2)}.
\end{equation*}
在动量空间，二点反项只能产生
\begin{equation*}
 \Sigma_{\rm ct}(p^2)
 =A_{\rm ct}+B_{\rm ct}(p^2-m_R^2c^2),
\end{equation*}
其中常数 $A_{\rm ct}$ 由质量反项控制，斜率 $B_{\rm ct}$ 由 $\delta Z_\varphi$ 控制。由于一圈蝌蚪满足
\begin{equation*}
 \frac{\dd\Sigma_{\rm tad}^{(1)}}{\dd p^2}=0,
\end{equation*}
一圈阶没有任何紫外斜率需要 $\delta Z_\varphi$ 抵消，因此
\begin{equation*}
 \delta Z_\varphi^{(1)}=0,
 \qquad
 Z_\varphi=1+O(g_4^2),
\end{equation*}
而常数紫外极点必须由 $\delta m^2=O(g_4)$ 吸收。这得到\eqnrefs{eq:phi4-one-loop-mass-vs-Z-dp52}。

最后检查壳上条件。精确重整化传播子定义为
\begin{equation*}
 \widetilde G_R(p)
 =\frac{\ii}
 {p^2-m_R^2c^2-\Sigma_R(p^2)+\ii0^+}.
\end{equation*}
要求物理极点仍位于 $p^2=m_R^2c^2$，分母在该点必须为零，因此
\begin{equation*}
 \Sigma_R(m_R^2c^2)=0.
\end{equation*}
在极点附近令 $p^2=m_R^2c^2+\delta$，Taylor 展开
\begin{align*}
 p^2-m_R^2c^2-\Sigma_R(p^2)
 &=\delta
 -\Sigma_R(m_R^2c^2)
 -\delta\,\Sigma_R'(m_R^2c^2)
 +O(\delta^2),\\
 &=\delta\,[1-\Sigma_R'(m_R^2c^2)]
 +O(\delta^2).
\end{align*}
若把重整化场定义成极点留数为 1，就必须再要求
\begin{equation*}
 \Sigma_R'(m_R^2c^2)=0.
\end{equation*}
这正是\eqnrefs{eq:phi4-onshell-ren-conditions-dp52}。因此质量反项与场强反项分别控制“极点位置”和“极点留数”；它们不是两个形式相似但物理意义不明的常数。
\end{exbox}

四维 \(\varphi^4\) 理论把“正则化—反项—尺度流”三步直接连在一起。一圈四点函数的三个通道都归结为同一个标量泡图。以 \(s\) 通道为例，先把需要计算的标量泡图明确写成
\begin{equation}
 {
 I(s)=\mu_{\rm DR}^{2\epsilon}
 \int\frac{\dd^d \ell}{(2\pi)^d}
 \frac{\ii}
 {\bigl(\ell^2-m_R^2c^2+\ii0^+\bigr)
  \bigl((\ell+p)^2-m_R^2c^2+\ii0^+\bigr)},
 \qquad p^2=s.}
 \label{eq:phi4-bubble-integral-r68}
\end{equation}
对\eqnrefs{eq:phi4-bubble-integral-r68} 做 Feynman 参数化和维数正则化后得到
\begin{equation}
 {
 I(s)
 =-\frac1{16\pi^2}
 \left[
 \frac1{\bar\epsilon}
 -\int_0^1\dd x\,
 \ln\frac{m_R^2c^2-x(1-x)s-\ii0^+}{\mu_{\rm DR}^2}
 \right]
 +\order(\epsilon),}
 \label{eq:phi4-bubble-result-dp10}
\end{equation}
其中 \(1/\bar\epsilon\equiv1/\epsilon-\gamma_E+\ln4\pi\)。三个通道具有相同的局域紫外极点，因此在 $d=4-2\epsilon$ 中，$\overline{\rm MS}$ 的一圈裸耦合可以写为
\begin{equation}
 {
 g_{4,B}
 =\mu^{2\epsilon}
 \left[
 g_4(\mu)
 +\frac{3}{32\pi^2}\frac{g_4^2(\mu)}{\bar\epsilon}
 +\order(g_4^3)
 \right].}
 \label{eq:phi4-bare-coupling-r68}
\end{equation}
裸参数不能依赖人为引入的减除尺度。对\eqnrefs{eq:phi4-bare-coupling-r68} 在固定 $g_{4,B}$ 下求尺度导数，得到
\begin{equation}
 {
 \beta_{g_4}(g_4)
 \equiv
 \mu\frac{\dd g_4}{\dd\mu}\Big|_{g_{4,B}}
 =\frac{3g_4^2}{16\pi^2}+\order(g_4^3).}
 \label{eq:phi4-beta-one-loop-dp10}
\end{equation}
一圈积分解为
\begin{equation}
 {
 \frac1{g_4(\mu)}
 =\frac1{g_4(\mu_0)}
 -\frac{3}{16\pi^2}\ln\frac\mu{\mu_0}.}
 \label{eq:phi4-running-one-loop-dp10}
\end{equation}

\begin{exbox}[一圈运行的数值尺度]
设某个参考尺度 $\mu_0$ 上测得 $g_4(\mu_0)=1$。先把分辨尺度提高三个数量级，考察一个仍然温和的尺度比
\begin{equation*}
 \frac\mu{\mu_0}=10^3,
 \qquad
 \ln\frac\mu{\mu_0}=6.9078.
\end{equation*}
求新尺度的耦合及大对数控制量，再计算同一一圈方程形式上的 Landau 极点。
\textbf{解。} 一圈解给
\begin{align*}
 \frac1{g_4(\mu)}
 &=1-\frac{3}{16\pi^2}(6.9078)\\
 &\simeq0.8688,
\end{align*}
所以
\begin{equation*}
 g_4(10^3\mu_0)\simeq1.151.
\end{equation*}
耦合本身只增加约 $15\%$，但固定阶展开中真正控制大对数的是
\begin{equation*}
 \epsilon_{\log}
 \equiv\frac{g_4}{16\pi^2}\left|\ln\frac\mu{\mu_0}\right|
 \simeq4.37\times10^{-2},
\end{equation*}
仍明显小于 1，因此一圈结果在这个例子里仍是受控的。

继续把一圈解的分母外推到零，可得到形式上的 Landau 极点
\begin{equation*}
 \frac{\mu_L}{\mu_0}
 =\exp\!\left(\frac{16\pi^2}{3g_4(\mu_0)}\right)
 \simeq7.3\times10^{22}
 \qquad[g_4(\mu_0)=1].
\end{equation*}
这个巨大数字不应被误读成“一圈理论精确预言了一个真实新粒子能标”。它只说明正的一圈 $\beta$ 函数会让耦合随紫外尺度增大，并且若坚持把同一低阶方程外推足够远，最终一定离开微扰控制区。重整化群的价值正是把这种尺度积累显式量化出来。
\end{exbox}

普通固定阶一圈截断要求 \(g_4/(16\pi^2)\ll1\)；跨越很大的尺度比时还需 \(|g_4\ln(\mu/\mu_0)|/(16\pi^2)\ll1\)，否则应使用重整化群方程重求和大对数，而不是继续展开\eqnrefs{eq:phi4-bubble-result-dp10} 的有限阶对数。

任意重整化 $n$ 点函数的尺度依赖应从整个重整化参数空间一次建立。若裸场与重整化场满足 $\varphi_B=Z_\varphi^{1/2}\varphi_R$，同一组外点的裸与重整化关联函数满足
\begin{equation}
 {
 G_B^{(n)}=Z_\varphi^{n/2}G_R^{(n)}.}
 \label{eq:bare-ren-green-relation-r68}
\end{equation}
设理论含一组无量纲耦合 $\{g_a(\mu)\}$ 和一组质量参数 $\{m_b(\mu)\}$。耦合和质量对尺度的对数导数分别描述参数怎样随分辨尺度流动；场归一化因子 $Z_\varphi$ 的尺度导数则描述量子涨落使场偏离经典工程尺度律的程度。后一个量通常称为场的\emph{反常维数}；“反常”在这里只表示额外的量子尺度修正，并不表示理论发生了量子对称性反常。具体定义为
\begin{equation}
{
 \beta_{g_a}\equiv\mu\frac{\dd g_a}{\dd\mu}\bigg|_{\rm bare},
 \qquad
 \beta_{m_b}\equiv\mu\frac{\dd m_b}{\dd\mu}\bigg|_{\rm bare},
 \qquad
 \gamma_\varphi\equiv\frac12\mu\frac{\dd\ln Z_\varphi}{\dd\mu}\bigg|_{\rm bare}.}
\label{eq:rg-general-flow-definitions-dp74}
\end{equation}
固定所有裸参数时，$G_B^{(n)}$ 不依赖人为减除尺度，因此得到一般尺度方程


\begin{proposition}[尺度不变裸量的微分方程]
设裸量与重整化量的关系可微，尺度微分保持裸参数固定，异常维数与 $\beta$ 函数采用正文定义。则有
\begin{equation}
 {
 \left[
 \mu\frac{\partial}{\partial\mu}
 +\sum_a\beta_{g_a}\frac{\partial}{\partial g_a}
 +\sum_b\beta_{m_b}\frac{\partial}{\partial m_b}
 +n\gamma_\varphi
 \right]G_R^{(n)}=0.}
\label{eq:callan-symanzik-general-dp10}
\end{equation}
\end{proposition}
\begin{proof}
固定所有裸参数，对正文\eqnrefs{eq:bare-ren-green-relation-r68} 求 $\mu$ 的全导数：
\begin{equation*}
 0=\mu\frac{\dd}{\dd\mu}
 \left[Z_\varphi^{n/2}G_R^{(n)}(\mu,\{g_a\},\{m_b\})\right].
\end{equation*}
除以 $Z_\varphi^{n/2}$ 后，场强因子给出
\begin{equation*}
 \mu\frac{\dd}{\dd\mu}\ln Z_\varphi^{n/2}
 =n\gamma_\varphi.
\end{equation*}
而 Green 函数的全导数沿全部重整化参数方向展开为
\begin{equation*}
 \mu\frac{\dd G_R^{(n)}}{\dd\mu}
 =\left(
 \mu\partial_\mu
 +\sum_a\mu\frac{\dd g_a}{\dd\mu}\partial_{g_a}
 +\sum_b\mu\frac{\dd m_b}{\dd\mu}\partial_{m_b}
 \right)G_R^{(n)}.
\end{equation*}
代入正文\eqnrefs{eq:rg-general-flow-definitions-dp74} 的 $\beta$ 函数定义便得到正文\eqnrefs{eq:callan-symanzik-general-dp10}。单耦合、单质量理论只是把上述求和各保留一项后的特例。
\end{proof}

这条方程描述重整化参数空间上的全尺度流。若只保留一个耦合 $g_{\rm R}$ 和一个质量 $m$，上式退化为
\[
 \left(
 \mu\partial_\mu+\beta_{g_{\rm R}}\partial_{g_{\rm R}}
 +\beta_m\partial_m+n\gamma_\varphi
 \right)G_R^{(n)}=0,
\]
这正是前面 $\varphi^4$ 一圈运行的参数特例。

局域复合算符具有另一种尺度结构。多个具有相同守恒量子数的算符在改变量子分辨尺度时可以彼此混合，因此它们的尺度修正不再由一个数描述，而由算符空间中的矩阵描述。这个矩阵就是复合算符的反常维数矩阵。若一组算符满足
\begin{equation}
{
 \mathcal O_i^B=Z_{ij}^{(\mathcal O)}\mathcal O_j^R,
 \qquad
 \gamma^{(\mathcal O)}
 =(Z^{(\mathcal O)})^{-1}\mu\frac{\dd Z^{(\mathcal O)}}{\dd\mu},}
\label{eq:operator-mixing-general-dp74}
\end{equation}
那么 $\gamma^{(\mathcal O)}$ 从一开始就是算符空间上的线性算符；其非对角元表示尺度流把一个基算符混入另一个基算符。只有在该空间恰为一维时，它才退化成单个反常维数。

\phantomsection\label{sec:qft-eft-unitarity-dp10}
重整化控制给定理论内部的紫外敏感性，却不要求低能实验显式解析所有高能自由度。若某些自由度的特征能量远高于实验能量，它们不能作为真实外态产生，但仍会通过虚过程改变低能振幅。

有效场论（EFT）只保留当前分辨率能够直接激发的轻自由度，并把重自由度的虚效应编码为由轻场构成的局域高维算符。若实验特征动量远小于重尺度，局域展开就按“实验动量/重尺度”的幂次有序；算符维数越高，通常受到越高次幂的抑制。EFT 的适用范围因此首先由这一无量纲展开参数控制，并还需满足散射矩阵幺正性等一致性界限。

采用 \(x^0=ct\)，定义几何拉氏量密度 \(\widehat\Lag=\Lag/(\hbar c)\)。同一个重尺度可等价写成逆长度、物理动量或能量：
\begin{equation}
 {
 \Lambda_*\ [\mathrm{m^{-1}}],
 \qquad
 p_*\equiv\hbar\Lambda_*,
 \qquad
 \Lambda_E\equiv cp_*=\hbar c\Lambda_*.}
 \label{eq:qft-eft-three-scales-dp42}
\end{equation}
若局域算符 \(\widehat{\mathcal O}_i\) 的逆长度工程维数为 \(\Delta_i>4\)，一般低能有效拉氏量为
\begin{equation}
 {
 \widehat\Lag_{\rm EFT}
 =\widehat\Lag_{\Delta\le4}
 +\sum_{\Delta_i>4}
 \frac{c_i}{\Lambda_*^{\Delta_i-4}}
 \widehat{\mathcal O}_i,}
 \label{eq:eft-operator-suppression-dp10}
\end{equation}
其中 \(c_i\) 无量纲。它们记录被消去高能自由度的短距离信息，而低能算符 \(\widehat{\mathcal O}_i\) 描述仍可分辨的自由度；乘在局域有效算符前的这些短距离系数称为\emph{Wilson 系数}。

Wilson 系数不能仅凭低能理论自身决定：应选择一组低能外态或 Green 函数，在包含重自由度的完整理论与有效理论中分别计算，并要求两种描述在共同适用的低能区域逐阶给出相同结果。用这种“比较同一低能物理并据此固定 Wilson 系数”的步骤称为\emph{匹配}。匹配解决的是高能信息怎样进入低能系数，而随后随重整化尺度改变这些系数则是另一个独立的重整化群问题。特征外部物理动量为 \(Q\) 时，维数 \(\Delta_i\) 的算符典型贡献按
\begin{equation}
 {
 \epsilon_{\rm EFT}^{(i)}
 \sim
 \left(\frac{Q}{p_*}\right)^{\Delta_i-4}
 \ll1}
 \label{eq:eft-power-counting-dp10}
\end{equation}
排序。若 \(Q\sim p_*\)，局域高维展开失去层级，必须恢复被消去的高能自由度或改用非局域描述。

有效场论把未解析的高能自由度编码为高维局域算符，并以 \(Q/p_*\) 排序；Feshbach--Schur、Schrieffer--Wolff 与 Green 张量投影则利用谱隙或预解式消去远失谐或连续自由度，控制参数通常是耦合与能隙之比以及预解式的能量依赖。二者都只保留给定能窗内可分辨的自由度，但展开参数不同。


有效场论截断还必须满足散射矩阵幺正性。对无自旋、单通道弹性 \(2\to2\) 散射，在质心系定义
\begin{equation}
 {
 \mathcal M(s,\theta)
 =16\pi\sum_{\ell=0}^{\infty}(2\ell+1)
 a_\ell(s)P_\ell(\cos\theta).}
 \label{eq:partial-wave-expansion-dp10}
\end{equation}
单弹性通道中 \(S_\ell=1+2\ii a_\ell\)，由 \(|S_\ell|=1\) 得到 Argand 圆，因此


\begin{proposition}[分波幺正约束]
设散射算符在完整开放通道空间幺正，分波归一化固定；单个弹性通道允许通过非弹性通道发生泄漏。则有
\begin{equation}
 {
 |a_\ell|\le1,
 \qquad
 |\Re a_\ell|\le\frac12.}
\label{eq:partial-wave-unitarity-bound-dp10}
\end{equation}
\end{proposition}
\begin{proof}
完整通道的幺正性给出弹性矩阵元 $|S_{\ell,\mathrm{el},\mathrm{el}}|\le1$。取 $S_{\ell,\mathrm{el},\mathrm{el}}=1+2\ii a_\ell$，则
\[
|1+2\ii a_\ell|^2\le1
\quad\Longleftrightarrow\quad
\left|a_\ell-\frac{\ii}{2}\right|\le\frac12.
\]
这个圆盘直接给出 $|a_\ell|\le1$、$|\Re a_\ell|\le1/2$，包含非弹性泄漏。没有泄漏的单通道弹性区间对应边界圆，下面计算其相移表示。

\begin{equation*}
 S_\ell=1+2\ii a_\ell,
 \qquad
 S_\ell S_\ell^*=1.
\end{equation*}
展开得到
\begin{align*}
 (1+2\ii a_\ell)(1-2\ii a_\ell^*)
 &=1+2\ii(a_\ell-a_\ell^*)+4|a_\ell|^2,\\
 &=1-4\Im a_\ell+4|a_\ell|^2.
\end{align*}
与 1 比较即有
\begin{equation*}
 \Im a_\ell=|a_\ell|^2.
\end{equation*}
写 \(a_\ell=x_\ell+\ii y_\ell\)，则
\begin{equation*}
 y_\ell=x_\ell^2+y_\ell^2,
\end{equation*}
配方得到
\begin{equation*}
 x_\ell^2+\left(y_\ell-\frac12\right)^2=\frac14.
\end{equation*}
因此精确弹性振幅位于以 \(\ii/2\) 为圆心、半径为 \(1/2\) 的 Argand 圆上。等价地，令
\begin{equation*}
 S_\ell=\ee^{2\ii\delta_\ell},
\end{equation*}
则
\begin{equation*}
 a_\ell
 =\frac{\ee^{2\ii\delta_\ell}-1}{2\ii}
 =\ee^{\ii\delta_\ell}\sin\delta_\ell,
\end{equation*}
从而
\begin{equation*}
 |a_\ell|\le1,
 \qquad
 |\Re a_\ell|\le\frac12.
\end{equation*}
树级振幅通常为实数，而精确弹性振幅必须位于上述 Argand 圆上。因此可采用保守的树级微扰一致性判据
\begin{equation*}
|\Re a_\ell^{\rm tree}|\le\frac12.
\end{equation*}
若树级结果越过该界，说明固定阶低能展开已经不能近似满足精确部分波幺正性，需要加入更高阶或新的有效自由度。
\end{proof}

存在非弹性通道时允许区域收缩到圆盘内部。若树级有效场论振幅已经增长到 \(|\Re a_\ell|\gtrsim1/2\)，失效的是固定阶低能截断；此时必须加入圈修正、重求和或新的高能自由度，而不能把精确理论解释为“不幺正”。

幂次计数判断潜在紫外行为；正则化定义发散积分；重整化用有限输入定义参数；重整化群组织减除尺度变化；有效场论按 \(Q/p_*\) 排列未解析高能结构；分波幺正性给出振幅的非微扰边界。QED、QCD 与电弱理论的规范群和场内容不同，但共享这套紫外组织原则。


这些结构可与 \cite{PeskinSchroeder1995,Schwartz2014QFT} 对照；全部尺度均按物理四动量与 SI 约定解释，并显式保留 \(c\) 与 \(\hbar\)。

% 多圈重整化、复合算符混合、算符乘积展开与有限温度边界条件
单圈原始发散图通常只有一个整体紫外发散，需要相应的局域反项吸收。进入多圈以后，同一图内部还会出现嵌套或重叠的发散真子图。若等整张图积分完成后再一次性减去极点，子图发散会被带入更高阶图中，无法由局域参数重定义逐阶吸收。局域性因此要求一种递归减除：先对所有真子图完成重整化，再处理缩约图剩余的整体发散。

对一张单粒子不可约图 $\Gamma$，记维数正则化后的裸振幅为 $I_\Gamma$。令线性算符 $K$ 只抽取所选减除方案中的局域紫外发散部分；例如在修正的最小减除方案（$\overline{\rm MS}$）中，$K$ 抽取所有 $1/\bar\epsilon^n$ 极点。对没有发散真子图的原始发散图，反项定义为
\begin{equation*}
 C(\Gamma)=-K I_\Gamma.
\end{equation*}
存在发散真子图时，先把所有真子图的局域反项插回原图，定义不完全 $R$ 运算
\begin{equation}
 \boxed{
 \bar R I_\Gamma
 =I_\Gamma
 +\sum_{\{\gamma_i\}}
 \left[\prod_i C(\gamma_i)\right]
 I_{\Gamma/\{\gamma_i\}}.}
 \label{eq:rbar-operation-dp68}
\end{equation}
其中求和遍历两两不重叠的紫外发散真子图集合，$\Gamma/\{\gamma_i\}$ 表示把每个 $\gamma_i$ 缩成一个局域顶角。完整减除的递归结构可表述如下。
\begin{proposition}[局域反项的递归减除]
固定正则化与线性紫外减除投影 $K$，并假设真子图递归减除后剩余的整体发散是 $K$ 所抽取的局域部分；红外奇异性已被另行调节。以两两不重叠的发散真子图集合定义 $\bar R$，其中各 $C(\gamma)$ 按较小子图递归定义。则整体反项和剩余振幅满足
\begin{equation}
 \boxed{
 C(\Gamma)=-K\bar RI_\Gamma,
 \qquad
 RI_\Gamma=(1-K)\bar RI_\Gamma}
 \label{eq:r-operation-dp11}
\end{equation}
\end{proposition}
\begin{proof}
先说明递归的代数结构。对没有发散真子图的图，$\bar RI_\gamma=I_\gamma$，于是 $C(\gamma)=-KI_\gamma$。假设所有较小真子图的反项已经定义，$\bar RI_\Gamma$ 是原图加上所有两两不重叠真子图集合的反项插入；把剩余整体反项定义为 $-K\bar RI_\Gamma$，立即有
\[
RI_\Gamma=\bar RI_\Gamma+C(\Gamma)=(1-K)\bar RI_\Gamma.
\]
图的真子图严格较小，所以归纳在有限步内结束。局域性与红外调节是假设；该代数恒等式本身不会证明这些分析条件。

逐层展开 $C(\gamma)=-K\bar RI_\gamma$ 时，顶层只选互不重叠的子图，下一层只选其内部的真子图。所得子图集合因而两两不重叠或互相包含；反过来，任一这样的森林由其最大真子图及各自的内部森林唯一确定。这给出递归项与兼容森林的一一对应，嵌套减除必须从内到外执行。重叠而互不包含的两个子图不能同时出现，分别属于不同森林。完整 $R$ 的最后一次整体减除还允许森林包含 $\Gamma$ 本身；只取真子图的森林描述的是 $\bar R$。

对只有一个发散真子图的两圈图，上述归纳的第一步可显式计算如下。

设两圈振幅为 $I_\Gamma$，其中 $\gamma$ 是一个对数发散的单圈真子图。对子图定义

\begin{equation*}
 C(\gamma)=-K I_\gamma.
\end{equation*}
在只有这一个发散真子图的情形，正文\eqnrefs{eq:rbar-operation-dp68} 化为
\begin{align*}
 \bar R I_\Gamma
 &=I_\Gamma+C(\gamma)I_{\Gamma/\gamma}\\
 &=I_\Gamma-(K I_\gamma)I_{\Gamma/\gamma}.
\end{align*}
第二项在每一个外层动量处抵消子图的大动量局域极点，因此 $\bar R I_\Gamma$ 只可能还保留整张图共同放大时的整体局域发散。正文\eqnrefs{eq:r-operation-dp11} 已把这一剩余整体发散的局域反项写成 $C(\Gamma)=-K\bar R I_\Gamma$。将该反项插回后便有
\begin{align*}
 R I_\Gamma
 &=\bar R I_\Gamma+C(\Gamma)\\
 &=(1-K)\bar R I_\Gamma,
\end{align*}
即正文\eqnrefs{eq:r-operation-dp11}。多个互不重叠或彼此嵌套的真子图只是在\eqnrefs{eq:rbar-operation-dp68} 中增加允许的子图集合，随后同样只对剩余整体发散作最后一次局域减除。
\end{proof}
第一式只抽取已经消除子图发散后剩余的整体局域极点；在上述分析假设下，第二式留下有限的重整化图。

嵌套与重叠子图可以用 Zimmermann 森林公式统一组织。用于完整 $R$ 的森林 $F$ 是一组彼此不重叠或互相包含的发散子图；若 $\Gamma$ 整体发散，也允许 $\Gamma$ 本身属于森林。以下明确取 BPHZ 的局域 Taylor 减除算符 $T_\gamma$，其取值阶数为相应表面发散度；它与前面最小减除中的极点投影 $K$ 是不同的方案实现。
\begin{corollary}[递归局域减除的森林组织]
假设 BPHZ 子图减除满足前述局域性条件，且所论外动量及红外调节允许相应 Taylor 减除。森林中的嵌套算符按从内到外的顺序作用，两两不重叠的减除可独立执行。则有限递归展开为
\begin{equation}
 \boxed{
 RI_\Gamma
 =\sum_{F\in\operatorname{Forests}(\Gamma)}
 \left(\prod_{\gamma\in F}(-T_\gamma)\right)I_\Gamma.}
 \label{eq:zimmermann-forest-dp11}
\end{equation}
\end{corollary}
\begin{proof}
对无发散真子图的图，森林仅含空集以及整体发散时的 $\{\Gamma\}$，求和给 $I_\Gamma-T_\Gamma I_\Gamma$。归纳到一般图：任一真子图森林的最大元素两两不重叠；每个最大元素内的子森林由较小图的归纳式给出。因此递归地选择最大元素、再选择其内部森林，恰好列举所有兼容的真子图森林，每项出现一次。最后整体减除把每个真森林项分别乘以 $I-T_\Gamma$，产生不含及包含 $\Gamma$ 的两项，得到所述完整森林和。互相重叠而不包含的子图不可能属于同一最大元素集合，因此不会被重复同时减除。此证明建立有限减除式的组合组织；紫外局域性与积分收敛仍依赖陈述中列明的分析假设。
\end{proof}
空森林给出原图；含一个子图的森林给一次减除；嵌套森林按从内到外的顺序作用；含整体 $\Gamma$ 的项完成最后的整体减除。定理的关键假设是紫外发散部分在外动量和轻质量上可作局域 Taylor 展开，因此 $T_\gamma$ 抽取的是有限阶局域多项式。这一局域性正是多圈反项仍能写回原拉格朗日量允许算符的原因。

重整化局域算符还需要处理重合点奇异性，因此场与耦合的重整化之后必须进一步引入复合算符混合。

分离时空点的 Green 函数完成重整化后，局域复合算符仍会产生额外的重合点奇异性。若一组算符 $\mathcal O_i$ 具有相同守恒量子数，重整化一般不是逐个乘一个常数，而是矩阵混合：
\begin{equation}
 \boxed{\mathcal O_{i,B}=Z_{ij}^{\mathcal O}\mathcal O_{j,R}.}
 \label{eq:composite-operator-mixing-dp11}
\end{equation}
定义反常维数矩阵
\begin{equation*}
 \gamma^{\mathcal O}
 \equiv (Z^{\mathcal O})^{-1}\mu\frac{\dd Z^{\mathcal O}}{\dd\mu},
\end{equation*}
即可得到
\begin{equation*}
 \mu\frac{\dd}{\dd\mu}\mathcal O_{i,R}
 =-\gamma_{ij}^{\mathcal O}\mathcal O_{j,R}.
\end{equation*}
若有效拉格朗日量含 $C_i(\mu)\mathcal O_i(\mu)$，物理作用量不能依赖任意重整化尺度，因此 Wilson 系数沿相反方向演化：
\begin{equation}
 \boxed{
 \mu\frac{\dd C_i}{\dd\mu}
 =C_j\gamma_{ji}^{\mathcal O}.}
 \label{eq:wilson-rg-dp11}
\end{equation}

\begin{proposition}[算符混合的对偶系数演化]
设重整化算符基封闭于线性混合，系数与算符的收缩表示同一个尺度不变的相互作用。则有以下结论（原式 \eqnrefs{eq:wilson-rg-dp11}）：
\begin{equation*}
 \boxed{
 \mu\frac{\dd C_i}{\dd\mu}
 =C_j\gamma_{ji}^{\mathcal O}.}
\end{equation*}
\end{proposition}
\begin{proof}
由

\begin{equation*}
 \mathcal O_B=Z^{\mathcal O}\mathcal O_R
\end{equation*}
并固定裸理论参数，
\begin{equation*}
 0=\mu\frac{\dd Z^{\mathcal O}}{\dd\mu}\mathcal O_R
 +Z^{\mathcal O}\mu\frac{\dd\mathcal O_R}{\dd\mu}.
\end{equation*}
左乘 $(Z^{\mathcal O})^{-1}$，并定义
\begin{equation*}
 \gamma^{\mathcal O}
 =(Z^{\mathcal O})^{-1}\mu\frac{\dd Z^{\mathcal O}}{\dd\mu},
\end{equation*}
得到
\begin{equation*}
 \mu\frac{\dd\mathcal O_{i,R}}{\dd\mu}
 =-\gamma_{ij}^{\mathcal O}\mathcal O_{j,R}.
\end{equation*}
有效作用量中的同一局域项写成 $C_i\mathcal O_i$。物理量不依赖任意重整化尺度，所以
\begin{align*}
 0
 &=\mu\frac{\dd}{\dd\mu}(C_i\mathcal O_i)\\
 &=\left(\mu\frac{\dd C_i}{\dd\mu}\right)\mathcal O_i
 -C_i\gamma_{ij}^{\mathcal O}\mathcal O_j.
\end{align*}
把第二项的哑指标交换到与第一项相同的算符基，逐个比较 $\mathcal O_i$ 的系数即得正文\eqnrefs{eq:wilson-rg-dp11}。
\end{proof}

多个局域算符之间的尺度演化都具有这一矩阵结构；具体物理系统只改变算符基与反常维数矩阵。

复合算符的尺度演化确定以后，短距离乘积可以继续按同一局域算符基展开，从而把短程系数与长程矩阵元分离。

当两个重整化局域算符的类空间隔 $x$ 远小于强子、介质或外态的长程尺度时，局域性允许把短距离乘积按局域算符基作渐近展开。算符乘积展开（OPE）的结构为
\begin{equation}
 \boxed{
 \mathcal O_A(x)\mathcal O_B(0)
 \sim\sum_i C_{AB}^{\ \ i}(x,\mu)\mathcal O_i(0,\mu),
 \qquad x\to0.}
 \label{eq:ope-general-dp11}
\end{equation}
Wilson 系数 $C_{AB}^{\ i}$ 编码距离尺度 $|x|$ 附近的短程物理，在硬尺度足够高时可用微扰理论计算。外态矩阵元 $\langle H|\mathcal O_i|H\rangle$ 则承载长程动力学。

OPE 按算符维数与对称性排序。截断误差由下一批被省略算符相对于保留项的幂次决定，而不能只看某一个 Wilson 系数是否数值很小。

\eqnrefs{eq:ope-general-dp11} 的两边必须具有相同的重整化尺度依赖。复合算符的混合因此强制 Wilson 系数满足耦合的矩阵重整化群方程；选择 $\mu$ 接近短距离硬尺度可以把大对数从固定阶系数中转移到尺度演化中。具体的短程--长程分离问题只需把相应算符基和外态代入这一一般结构。

自由场给出一个最直接的检查。Wick 定理给
\begin{equation*}
 \varphi(x)\varphi(0)
 =\mathcal C(x)+:\!\varphi(x)\varphi(0)\!:.
\end{equation*}
对正规序部分在 $x=0$ 附近作 Taylor 展开，
\begin{equation*}
 :\!\varphi(x)\varphi(0)\!:
 =:\!\varphi^2(0)\!:
 +x^\mu:\!(\partial_\mu\varphi)\varphi\!:(0)+\cdots.
\end{equation*}
单位算符的 Wilson 系数就是短距离奇异的传播子 $\mathcal C(x)$，而 $:\!\varphi^2\!:$ 等局域算符承载较软的场信息。相互作用理论中的 OPE 保留这一职责分离，但系数和算符都必须经过重整化并允许混合。

零温连续频率积分在热平衡态中还会因虚时间边界条件离散化。热密度算符由
\begin{equation*}
 \hat\rho_{\beta_T}=\frac{\ee^{-\beta_T\hat H}}{Z(\beta_T)},
 \qquad Z(\beta_T)=\Tr \ee^{-\beta_T\hat H},
 \qquad \beta_T=(k_BT)^{-1}
\end{equation*}
定义。把实时间解析延拓到虚时间 $t=-\ii\tau$ 后，迹的循环性给出 Kubo--Martin--Schwinger 边界条件：玻色场在 $\tau\to\tau+\beta_T\hbar$ 下周期，费米场反周期，
\begin{align}
 \phi(\tau+\beta_T\hbar)&=+\phi(\tau),\\
 \psi(\tau+\beta_T\hbar)&=-\psi(\tau).
 \label{eq:matsubara-boundary-dp11}
\end{align}
因此虚时 Fourier 模式只能取离散的 Matsubara 频率
\begin{equation}
 \boxed{
 \begin{aligned}
 \omega_n^{(B)}&=\frac{2\pi n}{\beta_T\hbar},\\
 \omega_n^{(F)}&=\frac{(2n+1)\pi}{\beta_T\hbar}.
 \end{aligned}}
 \label{eq:matsubara-frequencies-dp11}
\end{equation}

\begin{proposition}[周期与反周期 Fourier 频率]
在欧氏时间长度为 $\beta\hbar$ 的区间上，函数满足相应周期或反周期边界条件。则有以下结论（原式 \eqnrefs{eq:matsubara-frequencies-dp11}）：
\begin{equation*}
 \boxed{
 \begin{aligned}
 \omega_n^{(B)}&=\frac{2\pi n}{\beta_T\hbar},\\
 \omega_n^{(F)}&=\frac{(2n+1)\pi}{\beta_T\hbar}.
 \end{aligned}}
\end{equation*}
\end{proposition}
\begin{proof}
玻色 Fourier 模写成 $\ee^{-\ii\omega_n\tau}$。周期条件给

\begin{equation*}
 \ee^{-\ii\omega_n\beta_T\hbar}=1,
\end{equation*}
因此
\begin{equation*}
 \omega_n^{(B)}=\frac{2\pi n}{\beta_T\hbar},
 \qquad n\in\mathbb Z.
\end{equation*}
费米场满足反周期条件，于是
\begin{equation*}
 \ee^{-\ii\omega_n\beta_T\hbar}=-1,
\end{equation*}
从而
\begin{equation*}
 \omega_n^{(F)}=\frac{(2n+1)\pi}{\beta_T\hbar}.
\end{equation*}
两者合起来就是正文\eqnrefs{eq:matsubara-frequencies-dp11}。在区间 $0\le\tau<\beta_T\hbar$ 上，这些模式满足
\begin{equation*}
 \frac1{\beta_T\hbar}\int_0^{\beta_T\hbar}\dd\tau\,
 \ee^{\ii(\omega_n-\omega_m)\tau}=\delta_{nm},
\end{equation*}
因此零温连续 Fourier 测度相应变为正文\eqnrefs{eq:matsubara-loop-replacement-dp11} 的离散求和。
\end{proof}

零温圈积分中的连续频率测度相应替换为
\begin{equation}
 \boxed{
 \int\frac{\dd\omega}{2\pi}
 \longrightarrow
 \frac1{\beta_T\hbar}\sum_n.}
 \label{eq:matsubara-loop-replacement-dp11}
\end{equation}
例如自由玻色子的虚时传播子分母为 $(\omega_n^{(B)})^2+\omega_{\mathbf k}^2$。完成频率求和后出现 Bose--Einstein 分布；统计因子由热边界条件与传播子谱极点共同决定。



多圈 $R$ 运算要求紫外反项保持局域；复合算符重整化把重合点奇异性组织成混合矩阵；算符乘积展开再把短距离奇异性编码进 Wilson 系数。三者描述的是同一局域量子场论在不同尺度组织方式下的同一个紫外结构。


泛函行列式的固有时表示把一圈紫外问题转化为短固有时的局域渐近展开。对欧氏正谱二次算符 \(\Delta\)，先选取逆长度减除尺度 \(\mu_\kappa\,[\mathrm{m^{-1}}]\)。谱定理把标量恒等式 \(\ln(\lambda/\mu_\kappa^2)\) 提升为


\begin{proposition}[正算子行列式的固有时表示]
设欧氏二次算子为正自伴，先以有限模或热核截断定义相对行列式，再按正文做紫外减除；非正谱需另选解析延拓。则有
\begin{equation}
 {
 \Tr\ln\frac{\Delta}{\mu_\kappa^2}
 =-\int_0^\infty\frac{\dd s}{s}
 \Tr\!\left(\ee^{-s\Delta}-\ee^{-s\mu_\kappa^2}\right).}
\label{eq:proper-time-functional-dp11}
\end{equation}
\end{proposition}
\begin{proof}
正文\eqnrefs{eq:one_loop_effective_action_general} 表明一圈修正由二次涨落算符的 $\Tr\ln$ 决定。为把紫外问题改写成可按短距离展开的积分，先 Wick 转动到正谱欧氏算符 $\Delta$，再把标量对数恒等式逐个作用在 $\Delta$ 的本征值上。

对正数 $\lambda>0$ 定义
\begin{equation*}
 F(\lambda)
 =-\int_0^\infty\frac{\dd s}{s}
 \left(\ee^{-s\lambda}-\ee^{-s\mu_\kappa^2}\right).
\end{equation*}
对 $\lambda$ 求导，
\begin{equation*}
 \frac{\dd F}{\dd\lambda}
 =\int_0^\infty\dd s\,\ee^{-s\lambda}
 =\frac1\lambda,
\end{equation*}
并且 $F(\mu_\kappa^2)=0$，所以
\begin{equation*}
 F(\lambda)=\ln\frac{\lambda}{\mu_\kappa^2}.
\end{equation*}
若 $\Delta|n\rangle=\lambda_n|n\rangle$，谱定理给出
\begin{align*}
 \Tr\ln\frac{\Delta}{\mu_\kappa^2}
 &=\sum_n\ln\frac{\lambda_n}{\mu_\kappa^2}\\
 &=-\sum_n\int_0^\infty\frac{\dd s}{s}
 \left(\ee^{-s\lambda_n}-\ee^{-s\mu_\kappa^2}\right)\\
 &=-\int_0^\infty\frac{\dd s}{s}
 \Tr\!\left(\ee^{-s\Delta}-\ee^{-s\mu_\kappa^2}\right),
\end{align*}
这就是正文\eqnrefs{eq:proper-time-functional-dp11}。
\end{proof}

因为 \(\Delta\) 的本征值具有 \(\mathrm{m^{-2}}\) 量纲，固有时 \(s\) 具有 \(\mathrm{m^2}\) 量纲；大本征值对应 \(s\to0^+\)，所以紫外发散完全由热核的短时系数决定。若改用物理动量尺度，则 \(\mu_p=\hbar\mu_\kappa\)，不能把二者直接混入同一个无量纲比值。

在平直欧氏空间中，对 Laplace 型算符定义
\begin{equation}
 {
 \Delta=-\left(\nabla^2+E\right),
 \qquad
 \nabla_\mu=\partial_\mu+\omega_\mu,
 \qquad
 \Omega_{\mu\nu}=[\nabla_\mu,\nabla_\nu].}
 \label{eq:laplace-type-operator-dp11}
\end{equation}
其中 \(E\) 是纤维内部空间上的局域端同态，\(\omega_\mu\) 是联络，\(\Omega_{\mu\nu}\) 是其曲率。对角热核的短时展开为


\begin{proposition}[Laplace 型算子的局部热核系数]
设底空间为平直欧氏空间，算子为光滑 Laplace 型算子，采用正文的联络曲率与势项符号；在远离边界的内点，局部展开在 $s\downarrow0$ 时按渐近意义理解。下式成立；一般曲底空间还须加入 Riemann 曲率项：
\begin{equation}
 {
 K(x,x;s)
 =\frac1{(4\pi s)^{d/2}}
 \left[
 I+sE+s^2\left(
 \frac12E^2+\frac1{12}\Omega_{\mu\nu}\Omega^{\mu\nu}
 +\frac16\nabla^2E
 \right)+\order(s^3)
 \right].}
\label{eq:heat-kernel-a4-flat-dp11}
\end{equation}
\end{proposition}
\begin{proof}
在基点 $x$ 选择 Fock--Schwinger 规范

\begin{equation*}
 (y-x)^\mu\omega_\mu(y)=0,
\end{equation*}
则联络在 $x$ 附近为
\begin{equation*}
 \omega_\mu(y)
 =\frac12(y-x)^\nu\Omega_{\nu\mu}(x)
 +O((y-x)^2).
\end{equation*}
自由热核是
\begin{equation*}
 K_0(x,y;s)=\frac1{(4\pi s)^{d/2}}
 \exp\!\left[-\frac{(x-y)^2}{4s}\right].
\end{equation*}
令 $\xi=y-x$，其高斯矩为
\begin{align*}
 \langle\xi^\mu\xi^\nu\rangle_s
 &=2s\,\delta^{\mu\nu},\\
 \langle\xi^\mu\xi^\nu\xi^\rho\xi^\sigma\rangle_s
 &=4s^2(\delta^{\mu\nu}\delta^{\rho\sigma}
 +\delta^{\mu\rho}\delta^{\nu\sigma}
 +\delta^{\mu\sigma}\delta^{\nu\rho}).
\end{align*}
把 $E$ 与联络项作为固有时有序展开中的插入：一次 $E$ 插入给 $sE$，两次 $E$ 插入给 $s^2E^2/2$；两个线性联络插入结合四阶高斯矩给
\begin{equation*}
 \frac{s^2}{12}\Omega_{\mu\nu}\Omega^{\mu\nu},
\end{equation*}
而 $E(y)$ 的二阶协变 Taylor 项结合二阶高斯矩给
\begin{equation*}
 \frac{s^2}{6}\nabla^2E.
\end{equation*}
将这些项与自由热核前因子合并，就得到正文\eqnrefs{eq:heat-kernel-a4-flat-dp11}。
\end{proof}

无边界体积分中 \(\nabla^2E\) 是全导数，因此四维对数紫外极点真正依赖的局域不变量是 \(E^2\) 与 \(\Omega_{\mu\nu}\Omega^{\mu\nu}\)。

在 \(d=4-2\epsilon\) 中，\eqnrefs{eq:heat-kernel-a4-flat-dp11} 的 \(s^2\) 项与 \(\dd s/s\) 组合成 \(\dd s\,s^{-1+\epsilon}\)，从而直接产生 \(1/\epsilon\) 极点。对实玻色场的一圈高斯行列式，发散部分为
\begin{equation}
 {
 \left.\frac12\Tr\ln\Delta\right|_{\rm div}
 =-\frac1{2(4\pi)^2\bar\epsilon}
 \int\dd^4 x\,
 \tr\!\left[
 \frac12E^2+\frac1{12}\Omega_{\mu\nu}\Omega^{\mu\nu}
 \right].}
 \label{eq:heat-kernel-divergence-master-dp11}
\end{equation}
反交换标量变量和 Dirac 费米变量会因统计性质改变泛函行列式的整体符号与简并度，但允许出现的局域算符基仍由同一短时系数决定。对非阿贝尔规范场，规范场本身与规范固定 Jacobian 产生的 Grassmann 辅助场分别贡献各自的有限维指标迹；这些辅助变量的出现来自规范轨道换元的 Jacobian，而不是额外的物理渐近自由度。


% 多圈重整化、算符混合与虚时边界条件的集中推导。









\subsection{规范固定与量子规范冗余}
普通高斯场的传播子来自可逆二次核；规范势却含有不改变场强的冗余方向，所以其二次核有零模，不能直接求逆。必须在每条规范轨道上只计数一个代表，这就是规范固定。非阿贝尔情况下，这个变量约束产生依赖规范场的泛函 Jacobian；用一对反交换辅助场表示该 Jacobian 后，扩大场空间具有一个幂零对称变换。这个幂零结构把规范冗余转化为代数约束，并使物理态能够由相应上同调类选出。
% 一般规范场路径积分、Faddeev--Popov 规范固定与 BRST 结构
局域规范对称性的出发点可以先用一个简单事实理解：同一个物理内部状态，可以在每个时空点选择不同的内部基底来描述。若这个基底变化依赖于位置，普通导数会额外微分到基底本身，于是两个相邻点上的内部向量不能直接比较。为补偿这种“局域基底变化”而引入的规范势称为\emph{规范联络}；协变导数就是利用它规定相邻时空点之间怎样比较内部状态。电磁势 $A_\mu$ 对应阿贝尔 $U(1)$ 联络；当内部变换不再彼此交换时，同一局域比较原则给出非阿贝尔联络。

量子理论还遇到第二个问题：不同的联络函数可以只相差一次局域规范变换，却表示完全相同的物理状态。把“所有可能的场函数”整体看成一个巨大空间，称为\emph{场空间}；其中一个“点”不是普通坐标，而是一整套函数 $\mathscr C_\mu^a(x)$。从某个场构形出发施加所有可能的局域规范变换，会在场空间中扫出一族彼此物理等价的点，这一族称为\emph{规范轨道}。若路径积分把轨道上的每一点都积分一次，就会把同一个物理状态无限重复计数。

因此规范固定不是增加新的动力学，而是在每条规范轨道上选择一个代表。几何上，这相当于在场空间中选一张横截轨道的“切片”。换变量时还必须乘上 Jacobian；这里的\emph{泛函 Jacobian}只是普通多变量积分 Jacobian 行列式在无限多个场自由度上的对应物。它衡量规范参数发生微小变化时，规范条件函数改变多少。这个 Jacobian 可写成明确的泛函行列式；其标准名称是 Faddeev--Popov 行列式。

令紧致李群 $G$ 表示内部规范变换群，生成元与结构常数沿用第一章的 Lie 代数定义。所有联络构成的场空间记为 $\mathscr A$，一条规范轨道记为 $[\mathscr C]=\{\mathscr C^U\}$。作用量沿轨道不变，因此二次变分在轨道切向上有零模；未经规范固定时，传播子需要反演的二次核便不可逆。采用几何联络归一化
\begin{equation}
 \mathscr C_\mu=\mathscr C_\mu^aT^a,
 \qquad
 D_\mu=\partial_\mu+\ii g_{\rm G}\mathscr C_\mu,
 \label{eq:generic-gauge-cov-derivative-dp11}
\end{equation}
其中 $g_{\rm G}$ 是与一般规范群 $G$ 对应的无量纲规范耦合，$[\mathscr C_\mu]=\mathrm{m^{-1}}$。场强由协变导数的对易子定义，
\begin{align}
 [D_\mu,D_\nu]&=\ii g_{\rm G}\mathcal F_{\mu\nu},\\
 \mathcal F_{\mu\nu}^a
 &=\partial_\mu\mathscr C_\nu^a-\partial_\nu\mathscr C_\mu^a
 -g_{\rm G} f^{abc}\mathscr C_\mu^b\mathscr C_\nu^c,
 \label{eq:generic-field-strength-dp11}
\end{align}
对应的纯规范场拉格朗日密度为
\begin{equation}
 \Lag_{\rm YM}
 =-\frac{\hbar c}{4}\mathcal F_{\mu\nu}^a\mathcal F^{a\mu\nu}.
 \label{eq:generic-ym-lagrangian-dp11}
\end{equation}

一般局域规范变换首先作用在联络本身。与上面的协变导数约定相容的有限变换为
\begin{equation}
\boxed{
 \mathscr C_\mu'
 =U\mathscr C_\mu U^{-1}
 +\frac{\ii}{g_{\rm G}}(\partial_\mu U)U^{-1}.}
\label{eq:generic-finite-gauge-transform-dp68}
\end{equation}
现在才取无穷小形式 $U=I+\ii\vartheta^aT^a+\order(\vartheta^2)$。它给出规范轨道在场空间中的切向量
\begin{equation}
 \boxed{
 \delta\mathscr C_\mu^a
 =-\frac1{g_{\rm G}}(\mathcal D_\mu^{\rm adj})^{ab}\vartheta^b,
 \qquad
 (\mathcal D_\mu^{\rm adj})^{ab}
 =\delta^{ab}\partial_\mu-g_{\rm G}f^{acb}\mathscr C_\mu^c.}
 \label{eq:generic-infinitesimal-gauge-dp11}
\end{equation}

\begin{proposition}[有限规范变换的无穷小形式]
设群作用与联络随参数可微，取单位元附近的规范变换并保留参数一次项。则有以下结论（原式 \eqnrefs{eq:generic-infinitesimal-gauge-dp11}）：
\begin{equation*}
 \boxed{
 \delta\mathscr C_\mu^a
 =-\frac1{g_{\rm G}}(\mathcal D_\mu^{\rm adj})^{ab}\vartheta^b,
 \qquad
 (\mathcal D_\mu^{\rm adj})^{ab}
 =\delta^{ab}\partial_\mu-g_{\rm G}f^{acb}\mathscr C_\mu^c.}
\end{equation*}
\end{proposition}
\begin{proof}
正文\eqnrefs{eq:generic-finite-gauge-transform-dp68} 给出了联络在有限局域变换下的变化。把 $U$ 展开到规范参数 $\vartheta^a$ 的一阶，并使用 Lie 代数对易关系，就得到沿规范轨道的无穷小切向量~\eqnrefs{eq:generic-infinitesimal-gauge-dp11}。这一步只是在场空间中线性化有限变换。

从有限规范变换
\begin{equation*}
 \mathscr C_\mu'
 =U\mathscr C_\mu U^{-1}
 +\frac{\ii}{g_{\rm G}}(\partial_\mu U)U^{-1}
\end{equation*}
出发。令 $U=I+\ii\vartheta+\order(\vartheta^2)$、$U^{-1}=I-\ii\vartheta+\order(\vartheta^2)$，则
\begin{align*}
 U\mathscr C_\mu U^{-1}
 &=\mathscr C_\mu+\ii[\vartheta,\mathscr C_\mu]+\order(\vartheta^2),\\
 \frac{\ii}{g_{\rm G}}(\partial_\mu U)U^{-1}
 &=-\frac1{g_{\rm G}}\partial_\mu\vartheta+\order(\vartheta^2).
\end{align*}
因此
\begin{equation*}
 \delta\mathscr C_\mu
 =-\frac1{g_{\rm G}}\partial_\mu\vartheta+\ii[\vartheta,\mathscr C_\mu].
\end{equation*}
写成分量并使用 $[T^a,T^b]=\ii f^{abc}T^c$，得到
\begin{equation*}
 \delta\mathscr C_\mu^a
 =-\frac1{g_{\rm G}}\left(\partial_\mu\vartheta^a-g_{\rm G}f^{acb}\mathscr C_\mu^c\vartheta^b\right)
 =-\frac1{g_{\rm G}}(\mathcal D_\mu^{\rm adj})^{ab}\vartheta^b.
\end{equation*}
\end{proof}

同一规范轨道上的场构形给出相同物理状态，因此沿\eqnrefs{eq:generic-infinitesimal-gauge-dp11} 的切向方向积分只是在重复计算同一个物理点。

选取规范条件 $F^a[\mathscr C]=\omega^a$。它的作用只是从同一条规范轨道上选取一个代表。要把这一步放进路径积分，需要知道规范参数变化会把 $F$ 改变多少。把规范条件对无穷小规范参数的函数导数看成一个线性算符，它就是规范轨道坐标到规范条件坐标之间的 Jacobian 核；其泛函行列式给出积分测度的修正：
\begin{equation}
\boxed{
 \mathcal M_{\rm FP}^{ab}(x,y)
 =\left.\frac{\delta F^a[\mathscr C^\vartheta](x)}{\delta\vartheta^b(y)}\right|_{\vartheta=0},
 \qquad
 \Delta_{\rm FP}[\mathscr C]=\det\mathcal M_{\rm FP}.}
\label{eq:fp-jacobian-definition-dp68}
\end{equation}
算符 $\mathcal M_{\rm FP}$ 称为\emph{Faddeev--Popov 算符}，$\Delta_{\rm FP}$ 称为 Faddeev--Popov 行列式。若这个规范条件在所讨论的局部场空间区域内与每条规范轨道横截一次，函数空间的换元公式便给出
\begin{equation}
 \boxed{
 1=\Delta_{\rm FP}[\mathscr C]
 \int\mathcal D\vartheta\,
 \delta\!\left[F(\mathscr C^\vartheta)-\omega\right].}
 \label{eq:fp-identity-general-dp11}
\end{equation}

\begin{proposition}[局部 Faddeev--Popov 换元恒等式]
在有限维正则化且局部规范截面与每条轨道唯一相交的区域内，假设 FP 微分可逆；该结论不主张跨 Gribov 零点的全局有效性。则有以下结论（原式 \eqnrefs{eq:fp-identity-general-dp11}）：
\begin{equation*}
 \boxed{
 1=\Delta_{\rm FP}[\mathscr C]
 \int\mathcal D\vartheta\,
 \delta\!\left[F(\mathscr C^\vartheta)-\omega\right].}
\end{equation*}
\end{proposition}
\begin{proof}
规范固定的数学结构与普通多元积分中的换元完全相同，只是“坐标”变成了函数。先把规范场记为 \(\mathscr C_\mu^a(x)\)，无穷小规范参数记为 \(\vartheta^a(x)\)。沿规范轨道作变换
\[
\mathscr C_\mu
\longmapsto
\mathscr C_\mu^\vartheta .
\]
选定规范条件
\[
F^a[\mathscr C](x)=\omega^a(x).
\]
在一条给定规范轨道上，真正需要积分的是参数 \(\vartheta\)；因此定义映射
\[
\mathcal F_{\mathscr C}:
\vartheta
\longmapsto
F[\mathscr C^\vartheta]-\omega .
\]
若在所研究区域内每条规范轨道与切片 \(F=\omega\) 只有一个局部交点，则
\(\mathcal F_{\mathscr C}\) 在交点附近可逆。它的 Fr\'echet 导数为
\[
\left(
\mathcal M_{\rm FP}[\mathscr C]
\right)^{ab}(x,y)
\equiv
\left.
\frac{\delta F^a[\mathscr C^\vartheta](x)}
{\delta\vartheta^b(y)}
\right|_{\vartheta=0},
\]
即正文\eqnrefs{eq:fp-jacobian-definition-dp68}。

先回忆有限维公式。若 \(y=f(x)\) 且根 \(x_*\) 孤立，
\[
\delta^{(n)}[f(x)]
=
\frac{\delta^{(n)}(x-x_*)}
{|\det(\partial f/\partial x)_{x_*}|}.
\]
对 \(x\) 积分得到
\[
1
=
|\det J_f(x_*)|
\int\dd^n x\,\delta^{(n)}[f(x)].
\]
函数空间的局部换元完全平行：
\[
\mathcal D\!\left(F[\mathscr C^\vartheta]-\omega\right)
=
\det\mathcal M_{\rm FP}[\mathscr C^\vartheta]\,
\mathcal D\vartheta .
\]
于是只要切片在局部唯一，
\begin{align*}
\int\mathcal D\vartheta\,
\delta\!\left[F(\mathscr C^\vartheta)-\omega\right]
&=
\left[
\det\mathcal M_{\rm FP}[\mathscr C]
\right]^{-1},
\end{align*}
从而
\[
\boxed{
1
=
\Delta_{\rm FP}[\mathscr C]
\int\mathcal D\vartheta\,
\delta\!\left[
F(\mathscr C^\vartheta)-\omega
\right],
\qquad
\Delta_{\rm FP}[\mathscr C]
\equiv
\det\mathcal M_{\rm FP}[\mathscr C].
}
\]
这就是正文\eqnrefs{eq:fp-identity-general-dp11}。

为了看清这个抽象 Jacobian 的具体内容，取 Yang--Mills 场的无穷小规范变换
\[
\delta_\vartheta \mathscr C_\mu^a
=
(D_\mu\vartheta)^a
=
\partial_\mu\vartheta^a
+g_s f^{acb}\mathscr C_\mu^c\vartheta^b,
\]
并选 Lorenz 型规范条件
\[
F^a[\mathscr C]
=
\partial^\mu\mathscr C_\mu^a.
\]
则
\begin{align*}
\delta_\vartheta F^a(x)
&=
\partial^\mu
(D_\mu\vartheta)^a(x)\\
&=
\int\dd^4 y\,
\left[
\partial^\mu D_\mu^{ab}(x)
\delta^{(4)}(x-y)
\right]
\vartheta^b(y),
\end{align*}
所以
\[
\mathcal M_{\rm FP}^{ab}(x,y)
=
\partial^\mu D_\mu^{ab}(x)
\delta^{(4)}(x-y)
\]
到所选整体号规。Abelian \(U(1)\) 中 \(f^{abc}=0\)，于是
\[
\mathcal M_{\rm FP}
=
\partial^2
\]
与电磁场本身无关；其行列式只给一个可吸收的常数。非 Abelian 理论中 \(D_\mu\) 含规范场，因此
\(\det\mathcal M_{\rm FP}\) 是场依赖的，这正是后续 Faddeev--Popov ghost 与规范场耦合的来源。

若规范切片与同一轨道出现多个交点，上述“一个根对应一个 Jacobian”的局部换元不再全局成立；这就是 Gribov 问题出现的数学位置。Faddeev--Popov 恒等式因此是局部规范轨道坐标公式，而不是无条件的全局几何恒等式。
\end{proof}

这里 $\Delta_{\rm FP}$ 是从规范参数 $\vartheta$ 到规范条件 $F$ 的泛函雅可比，它只修正规范轨道上的积分测度。对协变规范 $F^a=\partial^\mu\mathscr C_\mu^a$，相应算符为
\begin{equation*}
 \mathcal M_{\rm FP}^{ab}
 =-\partial^\mu(\mathcal D_\mu^{\rm adj})^{ab}
\end{equation*}
（略去与场无关的整体常数）。对辅助函数 $\omega^a$ 作高斯平均后得到
\begin{equation}
 \boxed{
 \Lag_{\rm gf}
 =-\frac{\hbar c}{2\xi}
 (\partial^\mu\mathscr C_\mu^a)^2,}
 \label{eq:general-covariant-gf-dp11}
\end{equation}
其中 $\xi$ 是无量纲规范参数。

Faddeev--Popov 行列式仍是一个非局域的泛函行列式。前面已经见过 Grassmann 高斯积分会产生普通矩阵的行列式；把同一个代数恒等式应用到这里，可以用两组 Grassmann 奇的标量场 $c^a,\bar c^a$ 把该行列式写成局域作用量。它们称为\emph{鬼场}（ghost fields）。鬼场是表示规范轨道 Jacobian 的辅助 Grassmann 变量；BRST 物理态条件把它们排除在可探测的渐近态之外。由此得到鬼场拉格朗日量
\begin{equation}
 \boxed{
 \Lag_{\rm gh}
 =\hbar c\,\bar c^a
 \left[-\partial^\mu(\mathcal D_\mu^{\rm adj})^{ab}\right]c^b.}
 \label{eq:general-ghost-lagrangian-dp11}
\end{equation}
阿贝尔理论中 $f^{abc}=0$，$\mathcal M_{\rm FP}$ 与规范场无关，鬼场行列式可并入归一化；非阿贝尔理论中 $\mathcal M_{\rm FP}$ 含 $\mathscr C_\mu$，鬼场因而与规范场相互作用。鬼场不是物理渐近粒子，它们是规范轨道雅可比的局域场表示。

\begin{exbox}[自由 $U(1)$ 场中的规范零模]
一般规范轨道结构在自由电磁场中退化为线性零模。把 Fourier 波矢记为 $\kappa^\mu$，Maxwell 二次核为
\begin{equation}
 \boxed{
 \mathcal K_0^{\mu\nu}(\kappa)
 =-\kappa^2\eta^{\mu\nu}+\kappa^\mu\kappa^\nu,
 \qquad
 \mathcal K_0^{\mu\nu}\kappa_\nu=0.}
 \label{eq:u1-gauge-zero-mode-dp50}
\end{equation}
因此任意 $\mathscr C_\mu=\lambda\kappa_\mu$（$\lambda$ 为任意标量 Fourier 系数）都满足 $\mathcal K_0^{\mu\nu}\mathscr C_\nu=0$，它们正是规范轨道切向的零模。加入协变规范固定项 $-(\hbar c/2\xi)(\partial\!\cdot\!\mathscr C)^2$ 后，核变为
\begin{equation}
 \boxed{
 \mathcal K_\xi^{\mu\nu}
 =-\kappa^2\eta^{\mu\nu}
 +\left(1-\frac1\xi\right)\kappa^\mu\kappa^\nu,}
 \label{eq:u1-gauge-fixed-kernel-dp50}
\end{equation}
其横向与纵向本征值分别为 $-\kappa^2$ 与 $-\kappa^2/\xi$。这个线性阿贝尔例子展示一般规范轨道切片的最简单实现；Faddeev--Popov Jacobian 则由一般规范变换在规范条件上的泛函导数定义。

\medskip

\eqnrefs{eq:u1-gauge-zero-mode-dp50} 已经把问题具体化为“二次核沿 $\kappa^\mu$ 方向有零本征值”。这里加入协变规范固定项，并在横向、纵向投影上重新求二次核的本征值，从而得到可逆的\eqnrefs{eq:u1-gauge-fixed-kernel-dp50}。这一计算只涉及自由 $U(1)$ 场，还不需要鬼场或 BRST。

取 Fourier 约定 $\partial_\mu\to-\ii\kappa_\mu$。自由场强为
\begin{equation*}
 \mathcal F_{\mu\nu}
 =-\ii(\kappa_\mu\mathscr C_\nu-\kappa_\nu\mathscr C_\mu).
\end{equation*}
把它代入 $-\hbar c\,\mathcal F_{\mu\nu}\mathcal F^{\mu\nu}/4$，并把二次型写成
\begin{equation*}
 \frac{\hbar c}{2}\,
 \mathscr C_\mu(-\kappa)\mathcal K_0^{\mu\nu}(\kappa)\mathscr C_\nu(\kappa),
\end{equation*}
得到
\begin{equation*}
 \mathcal K_0^{\mu\nu}
 =-\kappa^2\eta^{\mu\nu}+\kappa^\mu\kappa^\nu.
\end{equation*}
右乘 $\kappa_\nu$ 后两项逐项抵消：
\begin{equation*}
 \mathcal K_0^{\mu\nu}\kappa_\nu
 =-\kappa^2\kappa^\mu+\kappa^\mu(\kappa^\nu\kappa_\nu)=0.
\end{equation*}
这就是规范变换 $\mathscr C_\mu\mapsto\mathscr C_\mu+\partial_\mu\chi$ 在 Fourier 空间中的切向方向。

协变规范固定项为
\begin{equation*}
 \Lag_{\rm gf}
 =-\frac{\hbar c}{2\xi}(\partial_\mu\mathscr C^\mu)^2.
\end{equation*}
它只修改纵向二次型，因此总核成为\eqnrefs{eq:u1-gauge-fixed-kernel-dp50}。定义横向、纵向投影
\begin{equation*}
 P_T^{\mu\nu}=\eta^{\mu\nu}-\frac{\kappa^\mu\kappa^\nu}{\kappa^2},
 \qquad
 P_L^{\mu\nu}=\frac{\kappa^\mu\kappa^\nu}{\kappa^2},
\end{equation*}
它们满足 $P_TP_L=0$、$P_T^2=P_T$、$P_L^2=P_L$。于是
\begin{equation*}
 \mathcal K_\xi
 =-\kappa^2 P_T-\frac{\kappa^2}{\xi}P_L.
\end{equation*}
因此除去 $\kappa^2=0$ 的物理传播极点以后，横向和纵向方向都已有非零本征值；规范固定正是通过给纯规范方向指定可逆的二次核来定义传播子。
\end{exbox}

规范固定除去路径积分中的重复计数以后，还需要保证不同规范参数给出相同的物理振幅。为此引入一种同时作用于规范场和鬼场的 Grassmann 奇变换；“Grassmann 奇”表示它会改变变量的玻色/费米奇偶性，并遵守带符号的 Leibniz 法则。这个变换称为 Becchi--Rouet--Stora--Tyutin（BRST）变换。其核心代数性质是幂零性 $s^2=0$：连续作用两次返回零。正是这一性质允许把规范冗余写成同调中的“恰当方向”，从而在量子态空间中识别物理等价类。

规范固定项可以写成一次 BRST 变换；为此引入 Nakanishi--Lautrup 辅助场 $B^a$。所谓\emph{辅助场}是指它没有独立传播的动能项，Euler--Lagrange 方程只是代数方程，因此可以在最后直接消去而不增加物理自由度。为与\eqnrefs{eq:generic-infinitesimal-gauge-dp11} 的规范变换约定一致，固定鬼场归一化为 $c^a=-\vartheta^a/g_{\rm G}$；BRST 微分 $s$ 定义为
\begin{align}
 s\mathscr C_\mu^a
 &=(\mathcal D_\mu^{\rm adj})^{ab}c^b,\\
 sc^a
 &=+\frac {g_{\rm G}}2 f^{abc}c^bc^c,\\
 s\bar c^a
 &=B^a,\\
 sB^a&=0.
 \label{eq:brst-transform-general-dp11}
\end{align}
连续作用两次时，真正需要检查的非平凡对象是鬼场和规范场；这不是两个彼此独立的结论，而是同一算符恒等式 $s^2=0$ 在两类非平凡场上的分量检验。分别有
\begin{align}
 s^2c^a&=0,
 \label{eq:brst-ghost-nilpotent-dp68}\\
 s^2\mathscr C_\mu^a&=0.
 \label{eq:brst-gauge-nilpotent-dp68}
\end{align}
第一式由结构常数的 Jacobi 恒等式与三个鬼场的反对称性给出；第二式还需要把 $\partial_\mu(sc)$ 与 $s\mathscr C_\mu$ 中的导数项逐项相消。再结合 $s\bar c^a=B^a$、$sB^a=0$，可汇总为
\begin{equation}
 \boxed{s^2=0}
 \label{eq:brst-nilpotent-dp11}
\end{equation}

\begin{proposition}[BRST 微分的幂零性]
设 Lie 代数结构常数满足 Jacobi 恒等式，ghost 为奇变量，$s$ 按分次 Leibniz 法则作用。则有以下结论（原式 \eqnrefs{eq:brst-nilpotent-dp11}）：
\begin{equation*}
 \boxed{s^2=0}
\end{equation*}
\end{proposition}
\begin{proof}
先检验鬼场。由
\begin{equation*}
 sc^a=\frac{g_{\rm G}}2f^{abc}c^bc^c
\end{equation*}
以及分次 Leibniz 法则
\begin{equation*}
 s(c^bc^c)=(sc^b)c^c-c^b(sc^c),
\end{equation*}
得到
\begin{align*}
 s^2c^a
 &=\frac{g_{\rm G}^2}{4}f^{abc}f^{bde}c^dc^ec^c
 -\frac{g_{\rm G}^2}{4}f^{abc}f^{cde}c^bc^dc^e.
\end{align*}
利用鬼场反对易性把三个鬼场统一排序，并对 $(b,d,e)$ 完全反对称化，
\begin{equation*}
 s^2c^a
 =\frac{g_{\rm G}^2}{12}
 \left(
 f^{abm}f^{mde}
 +f^{adm}f^{meb}
 +f^{aem}f^{mbd}
 \right)c^bc^dc^e=0,
\end{equation*}
其中最后一个等号使用 Lie 代数 Jacobi 恒等式。这就是正文\eqnrefs{eq:brst-ghost-nilpotent-dp68}。

再检验规范场。由
\begin{equation*}
 s\mathscr C_\mu^a
 =\partial_\mu c^a-g_{\rm G}f^{abc}\mathscr C_\mu^bc^c
\end{equation*}
再作用一次 $s$，
\begin{equation*}
 s^2\mathscr C_\mu^a
 =\partial_\mu(sc^a)
 -g_{\rm G}f^{abc}(s\mathscr C_\mu^b)c^c
 -g_{\rm G}f^{abc}\mathscr C_\mu^b(sc^c).
\end{equation*}
第一项化为
\begin{align*}
 \partial_\mu(sc^a)
 &=\frac{g_{\rm G}}2f^{ade}
 [ (\partial_\mu c^d)c^e+c^d\partial_\mu c^e]\\
 &=g_{\rm G}f^{ade}(\partial_\mu c^d)c^e,
\end{align*}
恰好抵消第二项中
$-g_{\rm G}f^{abc}(\partial_\mu c^b)c^c$ 的导数部分。剩余项先保持由 Leibniz 法则产生的系数，得到
\begin{align*}
 s^2\mathscr C_\mu^a
 &=g_{\rm G}^2f^{abc}f^{bde}\mathscr C_\mu^d c^e c^c
 -\frac{g_{\rm G}^2}{2}f^{abc}f^{cde}\mathscr C_\mu^b c^dc^e\\
 &=-\frac{g_{\rm G}^2}{2}\mathscr C_\mu^b
 \left(
 f^{abm}f^{mde}
 +f^{adm}f^{meb}
 +f^{aem}f^{mbd}
 \right)c^dc^e\\
 &=0,
\end{align*}
所以得到正文\eqnrefs{eq:brst-gauge-nilpotent-dp68}。此外
\begin{equation*}
 s^2\bar c^a=sB^a=0,
 \qquad
 s^2B^a=0,
\end{equation*}
故正文\eqnrefs{eq:brst-nilpotent-dp11} 对这组场变量全部成立。
\end{proof}

对上述全部场成立。幂零性不是记号便利，而是物理态上同调能够定义的必要条件。

规范固定项与鬼场项可以合并成一个 BRST 恰当项：
\begin{equation}
 \boxed{
 \Lag_{\rm gf+gh}
 =\hbar c\,s\left[
 \bar c^a\left(\partial^\mu\mathscr C_\mu^a+\frac\xi2B^a\right)
 \right].}
 \label{eq:brst-exact-gf-dp11}
\end{equation}
展开后为
\begin{equation}
 \frac{\Lag_{\rm gf+gh}}{\hbar c}
 =B^a\partial^\mu\mathscr C_\mu^a
 +\frac\xi2B^aB^a
 -\bar c^a\partial^\mu(\mathcal D_\mu^{\rm adj})^{ab}c^b,
 \label{eq:brst-gf-expanded-dp11}
\end{equation}

\begin{proposition}[BRST 恰当规范固定的展开]
设反 ghost 与辅助场的 BRST 变换采用正文约定，$s$ 为奇导子；按分次乘法展开规范固定费米子。则有以下结论（原式 \eqnrefs{eq:brst-gf-expanded-dp11}）：
\begin{equation*}
 \frac{\Lag_{\rm gf+gh}}{\hbar c}
 =B^a\partial^\mu\mathscr C_\mu^a
 +\frac\xi2B^aB^a
 -\bar c^a\partial^\mu(\mathcal D_\mu^{\rm adj})^{ab}c^b,
\end{equation*}
\end{proposition}
\begin{proof}
正文\eqnrefs{eq:brst-exact-gf-dp11} 把规范固定与鬼场项写成一个 BRST 恰当量。这里仅展开一次 BRST 微分，检查 $B$ 场、规范固定条件和 Faddeev--Popov 鬼算符的相对符号，从而得到正文\eqnrefs{eq:brst-gf-expanded-dp11}；随后消去代数辅助场 $B^a$ 的一步留在正文中完成。

令
\begin{equation*}
 \Psi_{\rm gf}
 =\bar c^a\left(\partial^\mu\mathscr C_\mu^a+\frac\xi2B^a\right).
\end{equation*}
因为 $\bar c$ 是 Grassmann 奇对象，分次 Leibniz 法则给
\begin{align*}
 s\Psi_{\rm gf}
 ={}&(s\bar c^a)
 \left(\partial\!\cdot\!\mathscr C^a+\frac\xi2B^a\right)
 -\bar c^a\left[
 \partial^\mu(s\mathscr C_\mu^a)+\frac\xi2sB^a
 \right].
\end{align*}
代入正文\eqnrefs{eq:brst-transform-general-dp11} 中的
$s\bar c^a=B^a$、$sB^a=0$ 与
$s\mathscr C_\mu^a=(\mathcal D_\mu^{\rm adj})^{ab}c^b$，得到
\begin{equation*}
 s\Psi_{\rm gf}
 =B^a\partial^\mu\mathscr C_\mu^a
 +\frac\xi2B^aB^a
 -\bar c^a\partial^\mu(\mathcal D_\mu^{\rm adj})^{ab}c^b.
\end{equation*}
乘回正文\eqnrefs{eq:brst-exact-gf-dp11} 中的整体因子 $\hbar c$，就是正文\eqnrefs{eq:brst-gf-expanded-dp11}。
\end{proof}

消去没有动力学导数的 $B^a$ 就恢复\eqnrefs{eq:general-covariant-gf-dp11,eq:general-ghost-lagrangian-dp11}。因此协变规范固定不是与规范理论并列的另一套动力学，而是同一规范轨道空间的一种坐标选择。

协变量子化为了保持 Lorentz 协变性，会暂时把横向、纵向、类时规范激发和鬼场都放进计算用的 Fock 空间。因此这个空间包含冗余方向，不能直接等同于物理 Hilbert 空间。

BRST 荷 $Q_B$ 在态空间上实现幂零关系 $Q_B^2=0$。物理态取为 BRST 闭合态，而相差一个 BRST 恰当态的两个闭合态代表同一物理状态：
\begin{equation*}
 Q_B|\mathrm{phys}\rangle=0,
 \qquad
 |\mathrm{phys}\rangle\sim|\mathrm{phys}\rangle+Q_B|\chi\rangle.
\end{equation*}
其中 $Q_B|\chi\rangle$ 只改变规范代表而不改变物理矩阵元。于是物理态空间是核与像之商 $\ker Q_B/\operatorname{im}Q_B$；这正是 BRST 上同调的定义。

对生成泛函作无穷小 BRST 换元时，积分测度与规范固定后的作用量保持不变。对单粒子不可约有效作用量 $\Gamma_{\rm eff}$，这一事实写成 Slavnov--Taylor 泛函恒等式
\begin{equation}
 \boxed{\mathcal S(\Gamma_{\rm eff})=0.}
 \label{eq:st-functional-general-dp11}
\end{equation}
对它作不同阶泛函微分可得到各个顶角函数之间的 Slavnov--Taylor 恒等式。阿贝尔 $U(1)$ 中鬼场脱耦，它退化为 Ward--Takahashi 恒等式；非阿贝尔理论中则同时约束物质场、规范场和鬼场顶角，以及波函数与耦合的重整化常数。

量子化还可能暴露另一类更根本的问题。一个连续变换即使保持经典作用量不变，调节后的函数积分测度或量子有效作用量也未必能同时保持同一对称性。若这种破坏不能通过允许的局域反项消去，就说这个经典对称性在量子理论中具有\emph{量子反常}（quantum anomaly）。对整体对称性，反常可以成为真实的量子效应；对必须用来消除非物理自由度的局域规范对称性，未抵消的规范反常会破坏物理态空间的一致性，因此可接受的基本规范理论必须使所有规范反常相互抵消。轴流反常与手征规范反常的差别，正来自这一一般量子测度效应作用在整体对称性或局域规范对称性上的不同物理角色。

Faddeev--Popov 行列式解决“规范轨道重复积分”的测度问题；鬼场把这个泛函行列式局域化；BRST 幂零性进一步保证规范固定后的扩大场空间能够通过上同调投影回物理态空间。三者属于同一规范冗余问题的不同层次，而不是彼此独立的附加规则。






\subsection{自发对称破缺}
一个对称势能可以拥有一整族等能最低点。理论方程仍保持对称，但一旦选择其中一个最低点作为真空，沿最低点族切向移动就可能没有恢复力。对真正作用在物理态上的整体连续对称性，这种平坦方向对应低能无隙激发；对局域规范冗余，同样的场方向会与规范场的纵向分量重新组合。为了区分这两种机制，先用势能 Hessian 判断真空切向方向是否具有恢复力，再用规范质量矩阵判断哪些方向被重新组织进有质量规范场。
对称的运动方程并不保证最低能量态本身也具有同样的对称性。最简单的情形是势能最低点不是一个点，而是一整圈等能真空：理论允许沿这条圆周变换而不改变能量，但一旦实际选定其中一个真空，围绕该真空的“径向”与“沿圆周”涨落就不再等价。沿圆周移动没有恢复力，因此会出现一个无能隙的低能模。连续整体对称性自发破缺时出现这种无隙模的结论通常称为 Goldstone 定理；只从真空几何和守恒荷把这一结论推出来。

这里还必须区分两种性质不同的变换。若同一个变换参数在整个时空中取同一常数，它连接不同的物理态，是整体对称性；若参数可以逐点任意改变而只是在描述同一物理配置，它属于规范冗余。前一种情形的无能隙模是真实激发，后一种情形则会在规范场存在时与纵向规范自由度重新组合。整体 \(U(1)\) 复标量场已经足以显示“简并真空为何产生无能隙方向”；把同一对称性局域化以后，角向涨落才与纵向规范自由度重新组合。

采用的 SI 场归一化
\begin{equation*}
 [\varphi_{\rm ssb}]=\sqrt{\mathrm{J/m}},
\end{equation*}
使 \((\partial_\mu\varphi_{\rm ssb})^*(\partial^\mu\varphi_{\rm ssb})\) 具有 \(\mathrm{J/m^3}\) 量纲。取
\begin{equation}
 \Lag
 =(\partial_\mu\varphi_{\rm ssb})^*(\partial^\mu\varphi_{\rm ssb})-V(\varphi_{\rm ssb}),
 \qquad
 V(\varphi_{\rm ssb})
 =-\mu_{\rm ssb}^2\varphi_{\rm ssb}^*\varphi_{\rm ssb}
 +\frac{\lambda}{\hbar c}(\varphi_{\rm ssb}^*\varphi_{\rm ssb})^2,
 \label{eq:ssb-complex-scalar-lag-dp11}
\end{equation}
其中 \(\mu_{\rm ssb}\,[\mathrm{m^{-1}}]\)、\(\lambda>0\) 无量纲。势能只依赖 \(\varphi_{\rm ssb}^*\varphi_{\rm ssb}\)，故在 \(\varphi_{\rm ssb}\to \ee^{\ii\alpha}\varphi_{\rm ssb}\) 下不变。

非平凡真空由势能极小条件确定。定义实真空尺度 \(v\) 使
\begin{equation}
 \boxed{
 \langle\varphi_{\rm ssb}\rangle=\frac{v}{\sqrt2},
 \qquad
 v^2=\frac{\mu_{\rm ssb}^2\hbar c}{\lambda}.}
 \label{eq:ssb-vev-general-dp11}
\end{equation}
所有 \((v/\sqrt2)\ee^{\ii\theta}\) 具有相同能量；选择其中一个真空破坏整体 \(U(1)\)，但作用量仍保持该对称性。

在所选真空附近写成 Cartesian 涨落
\begin{equation}
 \varphi_{\rm ssb}(x)=\frac1{\sqrt2}\left[v+h(x)+\ii\pi(x)\right].
 \label{eq:ssb-cartesian-fields-dp11}
\end{equation}
势能 Hessian 的二次项为
\begin{equation}
 V^{(2)}=\mu_{\rm ssb}^2h^2+0\times\pi^2,
 \label{eq:ssb-quadratic-potential-dp11}
\end{equation}
与实标量质量项 \(\tfrac12(m^2c^2/\hbar^2)\phi^2\) 比较可得
\begin{equation}
 \boxed{
 m_h^2=\frac{2\hbar^2\mu_{\rm ssb}^2}{c^2},
 \qquad
 m_\pi=0.}
 \label{eq:ssb-higgs-goldstone-masses-dp11}
\end{equation}
\(h\) 沿真空流形的径向振动，\(\pi\) 沿连续简并真空的切向振动。质量为零不是参数偶合，而是势能沿对称轨道恒定所强制的 Hessian 零本征值。

上面的 Hessian 计算依赖具体的四次势，但“沿被破缺对称方向没有能隙”并不依赖这个势的特殊形式。对任意连续整体对称性，Noether 流 \(j_a^\mu\) 定义守恒荷；若这个荷作用在某个局域场上能把真空推向一个不同真空，就有
\begin{equation}
\boxed{
 Q_a=\int\dd^3 x\,j_a^0(x),
 \qquad
 \frac{\dd Q_a}{\dd t}=0,
 \qquad
 \langle\Omega|[Q_a,\phi_i(0)]|\Omega\rangle\neq0.}
\label{eq:goldstone-broken-charge-dp68}
\end{equation}
最后一个非零对易子就是“真空没有保持该生成元”的可计算判据。无限体积中应把荷理解为局域对易子的空间截断极限；其谱结论与单粒子解释需要区分。
\begin{theorem}[守恒局域荷的无能隙判据]
设物理 Hilbert 空间的联合能量--动量谱位于正向光锥，真空平移不变且在所选纯真空扇区内是唯一的零能平移不变态。守恒局域流的空间截断荷 $Q_{a,R}$ 与局域场的真空对易子存在弱分布极限；流的局域性及空间边界通量条件保证该极限为不随时间变化的非零常数。假设完整态谱分解与这些涂抹及极限相容，则流与场共同耦合的非真空谱不能在零能、零动量附近有严格正能隙。若该谱支撑进一步由稳定单粒子分支实现，才可记作
\begin{equation}
\boxed{
 \exists\,|G_a(\mathbf p)\rangle:
 \qquad
 \lim_{\mathbf p\to0}E_{G_a}(\mathbf p)=0.}
\label{eq:goldstone-gapless-spectrum-dp68}
\end{equation}
\end{theorem}
\begin{proof}
在无限体积的破缺真空中，不能预先断言 $Q_a$ 是作用于真空的良定义全局算符。以下 $Q_a$ 及空间积分统一指 $Q_{a,R}$ 的局域对易子弱极限；不先对未涂抹的平面波交换无限体积极限。谱分解计算写为

设连续整体对称性有守恒流 \(j_a^\mu\) 与荷

\begin{equation*}
 Q_a=\int\dd^3 x\,j_a^0(x),
 \qquad
 \frac{\dd Q_a}{\dd t}=0.
\end{equation*}
若某个局域场满足
\begin{equation*}
 \langle\Omega|[Q_a,\phi_i(0)]|\Omega\rangle\neq0,
\end{equation*}
则把 \(Q_a\) 写成空间积分并在对易子中插入完整能量--动量本征态，得到
\begin{align*}
 \langle[Q_a,\phi_i(0)]\rangle
 =\int\dd^3 x\sum_n
 \bigl[&\langle\Omega|j_a^0(x)|n\rangle
 \langle n|\phi_i(0)|\Omega\rangle\\
 &-\langle\Omega|\phi_i(0)|n\rangle
 \langle n|j_a^0(x)|\Omega\rangle\bigr].
\end{align*}
平移不变性给
\begin{equation*}
 \langle\Omega|j_a^0(x)|n\rangle
 =\ee^{-\ii p_n\cdot x/\hbar}
 \langle\Omega|j_a^0(0)|n\rangle.
\end{equation*}
空间积分产生 \((2\pi\hbar)^3\delta^{(3)}(\mathbf p_n)\)，故只有零三动量态贡献。左边又因 \(Q_a\) 守恒与真空定态性而与时间无关；若所有零三动量中间态都有严格正的能隙，则右边只能含非零频率相位，不能维持非零常数。因此谱必须在 \(\mathbf p\to0\) 时具有 \(E\to0\) 的支撑，这给出 Goldstone 无隙谱判据；稳定单粒子解释仍需陈述中的附加条件。

这个谱论证还需排除把非零常数归因于真空或额外零能真空的可能性：真空中间态在两项对易子中相消，而所选纯真空扇区不含额外零能平移不变向量。更精确地，若存在统一正谱隙 $\Delta>0$，则两个非真空谱项的时间 Fourier 支撑分别落在 $\omega\ge\Delta/\hbar$ 和 $\omega\le-\Delta/\hbar$。弱极限相容性使它们的和仍在零频邻域为零；非零常数的 Fourier 变换却是非零的 $\delta(\omega)$，矛盾。空间荷截断的 Fourier 核集中到零三动量，所以这一障碍来自低动量的低能谱。于是结论首先是无隙谱支撑；只有再假设稳定单粒子分支，才得到陈述框中的 $|G_a(\mathbf p)\rangle$ 表示。

\end{proof}
Goldstone 判据所需的关键输入是物理局域场代数中的整体对称流、守恒局域荷的弱对易子极限及非零破缺序参量。局域规范变换属于描述冗余，不满足这一前提，必须另行分析其约束与物理自由度。

局域规范对称性的质量结构直接由一般群表示和标量真空的稳定子确定。设紧致规范群的一个简单因子为 $G$，标量多重态 $\varphi_{\rm ssb}$ 属于酉表示 $R$，其厄米生成元为 $T_R^a$，并定义
\begin{equation}
 D_\mu\varphi_{\rm ssb}
 =\left(\partial_\mu+\ii g_{\rm G}\mathscr C_\mu^aT_R^a\right)\varphi_{\rm ssb}.
 \label{eq:general-ssb-covariant-derivative-dp73}
\end{equation}
若势能的平稳真空为 $\langle\varphi_{\rm ssb}\rangle=v_R$，则在真空上
\begin{equation*}
 D_\mu v_R=\ii g_{\rm G}\mathscr C_\mu^aT_R^av_R.
\end{equation*}
因此标量动能的规范场二次项为
\begin{align*}
 (D_\mu v_R)^\dagger(D^\mu v_R)
 &=g_{\rm G}^2\mathscr C_\mu^a\mathscr C^{b\mu}
 v_R^\dagger T_R^aT_R^bv_R,\\
 &=\frac{g_{\rm G}^2}{2}\mathscr C_\mu^a\mathscr C^{b\mu}
 v_R^\dagger\{T_R^a,T_R^b\}v_R,
\end{align*}
其中第二行只用了 $\mathscr C_\mu^a\mathscr C^{b\mu}$ 对 $a,b$ 的对称性。与几何规范场归一化下的质量项
\[
\frac{\hbar c}{2}\mathscr C_\mu^a(M_g^2)_{ab}\mathscr C^{b\mu}
\]
比较，得到一般规范玻色子几何质量平方矩阵
\begin{equation}
 \boxed{
 (M_g^2)_{ab}
 =\frac{g_{\rm G}^2}{\hbar c}
 v_R^\dagger\{T_R^a,T_R^b\}v_R.}
 \label{eq:general-gauge-boson-mass-matrix-dp11}
\end{equation}
其本征值满足 $\lambda_r(M_g^2)=(m_rc/\hbar)^2$。更直接地，对任意实伴随向量 $x_a$，
\begin{equation}
 \boxed{
 x_a(M_g^2)_{ab}x_b
 =\frac{2g_{\rm G}^2}{\hbar c}
 \bigl\|x_aT_R^av_R\bigr\|^2\ge0.}
 \label{eq:gauge-mass-positive-quadratic-form-dp68}
\end{equation}

\begin{proposition}[规范质量矩阵的 Gram 正性]
设标量内积正定，质量矩阵由真空上的生成元作用与其厄米共轭构造。则有以下结论（原式 \eqnrefs{eq:gauge-mass-positive-quadratic-form-dp68}）：
\begin{equation*}
 \boxed{
 x_a(M_g^2)_{ab}x_b
 =\frac{2g_{\rm G}^2}{\hbar c}
 \bigl\|x_aT_R^av_R\bigr\|^2\ge0.}
\end{equation*}
\end{proposition}
\begin{proof}
正文\eqnrefs{eq:general-gauge-boson-mass-matrix-dp11} 给出生成元空间中的质量矩阵。要判断它是否可能产生负质量平方，不需要先对角化；只需把它与任意实向量 \(x_a\) 收缩，并把结果改写成真空向量的范数平方。


定义
\begin{equation*}
 (M_g^2)_{ab}
 =\frac{g_{\rm G}^2}{\hbar c}
 v_R^\dagger\{T^a,T^b\}v_R.
\end{equation*}
对任意实伴随空间向量 $x_a$，令
\begin{equation*}
 X\equiv x_aT^a.
\end{equation*}
若采用厄米生成元 $T^{a\dagger}=T^a$，则 $X^\dagger=X$。质量矩阵对应的二次型逐步化为
\begin{align*}
 x_a(M_g^2)_{ab}x_b
 &=\frac{g_{\rm G}^2}{\hbar c}
 v_R^\dagger x_ax_b\{T^a,T^b\}v_R,\\
 &=\frac{g_{\rm G}^2}{\hbar c}
 v_R^\dagger\{X,X\}v_R,\\
 &=\frac{2g_{\rm G}^2}{\hbar c}v_R^\dagger X^2v_R,\\
 &=\frac{2g_{\rm G}^2}{\hbar c}(Xv_R)^\dagger(Xv_R),\\
 &=\frac{2g_{\rm G}^2}{\hbar c}\|Xv_R\|^2\ge0.
\end{align*}
因此 $M_g^2$ 是实对称正半定矩阵，树级谱不会从这一机制产生负质量平方的矢量模。

等号成立当且仅当
\begin{equation*}
 Xv_R=0.
\end{equation*}
于是质量矩阵的零空间恰由所有湮灭真空的 Lie 代数方向组成，也就是未破缺子代数。反之，若 $Xv_R\neq0$，对应二次型严格为正，该生成元方向属于破缺子空间，并为相应规范玻色子提供非零质量。破缺方向上的 Goldstone 切向自由度与规范场纵向偏振重组为有质量矢量表示；这一线性代数结论不依赖某个具体规范群的显式矩阵对角化。
\end{proof}

因此 $M_g^2$ 正半定，而零空间恰好由满足 $x_aT_R^av_R=0$ 的 Lie 代数组合组成；这些方向保持真空不变，构成未破缺子代数。非零本征方向对应被真空选择破缺的生成元，并产生有质量规范玻色子。这个结论完全由一般表示和真空决定。

现在才取单生成元的 $U(1)$ 特例，以显式看清自由度如何重组。令标量的所选 $U(1)$ 生成元本征值吸收到 $g_{\rm G}$ 中，采用几何规范联络 $\mathscr C_\mu\,[\mathrm{m^{-1}}]$，
\begin{equation*}
 D_\mu=\partial_\mu+\ii g_{\rm G}\mathscr C_\mu,
 \qquad
 \mathcal F_{\mu\nu}
 =\partial_\mu\mathscr C_\nu-\partial_\nu\mathscr C_\mu.
\end{equation*}
带有前述自发破缺复标量的这个阿贝尔规范模型通常称为 Abelian Higgs 模型，其拉格朗日量为
\begin{equation}
 \Lag_{\rm AH}
 =(D_\mu\varphi_{\rm ssb})^*(D^\mu\varphi_{\rm ssb})-V(\varphi_{\rm ssb})
 -\frac{\hbar c}{4}\mathcal F_{\mu\nu}\mathcal F^{\mu\nu}.
 \label{eq:abelian-higgs-lag-dp11}
\end{equation}
用极坐标参数化
\begin{equation}
 \varphi_{\rm ssb}(x)=\frac{v+h(x)}{\sqrt2}
 \exp\!\left[\frac{\ii\pi(x)}{v}\right].
 \label{eq:abelian-higgs-polar-dp11}
\end{equation}
局域变换 $\varphi_{\rm ssb}\to \ee^{\ii\alpha(x)}\varphi_{\rm ssb}$ 对应
\begin{equation*}
 \pi\to\pi+v\alpha,
 \qquad
 \mathscr C_\mu\to\mathscr C_\mu-\frac1{g_{\rm G}}\partial_\mu\alpha,
\end{equation*}
所以组合
\begin{equation}
 \boxed{
 \mathcal B_\mu
 \equiv\mathscr C_\mu+\frac1{g_{\rm G}v}\partial_\mu\pi}
 \label{eq:gauge-invariant-massive-vector-dp11}
\end{equation}
在局域变换下不变。标量动能变为
\begin{equation}
 (D_\mu\varphi_{\rm ssb})^*(D^\mu\varphi_{\rm ssb})
 =\frac12(\partial h)^2
 +\frac{g_{\rm G}^2}{2}(v+h)^2\mathcal B_\mu\mathcal B^\mu.
 \label{eq:abelian-higgs-kinetic-expanded-dp11}
\end{equation}
其中二次项正是一般质量矩阵~\eqnrefs{eq:general-gauge-boson-mass-matrix-dp11} 在一维表示、$v_R=v/\sqrt2$ 时的特例。与
\[
 \frac{\hbar c}{2}
 \left(\frac{m_Ac}{\hbar}\right)^2
 \mathcal B_\mu\mathcal B^\mu
\]
比较得到
\begin{equation}
 \boxed{m_Ac^2=g_{\rm G}\sqrt{\hbar c}\,v.}
 \label{eq:abelian-higgs-vector-mass-dp11}
\end{equation}
若定义能量形式真空尺度 $v_E\equiv\sqrt{\hbar c}\,v$，则 $m_Ac^2=g_{\rm G}v_E$。破缺前的两个实标量自由度与两个无质量规范横向偏振总数为 4；破缺后重组为一个径向标量和三个有质量矢量偏振，自由度总数不变。

微扰计算若保留 $\mathscr C_\mu$ 与 $\pi$ 为独立场，需要选择能消去二次混合的规范固定。Cartesian 参数化产生
\begin{equation}
 \Lag^{(2)}\supset g_{\rm G}v\,\mathscr C^\mu\partial_\mu\pi.
 \label{eq:gauge-goldstone-mixing-dp11}
\end{equation}
定义
\begin{equation}
 \boxed{
 F=\partial^\mu\mathscr C_\mu
 -\xi\frac{g_{\rm G}v}{\hbar c}\,\pi,}
 \label{eq:rxi-gauge-functional-abelian-dp42}
\end{equation}
并加入
\begin{equation}
 \Lag_{\rm gf}
 =-\frac{\hbar c}{2\xi}F^2.
 \label{eq:rxi-abelian-gf-dp11}
\end{equation}
这种一参数规范选择通常称为 $R_\xi$ 规范。其交叉项经分部积分后抵消\eqnrefs{eq:gauge-goldstone-mixing-dp11}，剩余 $\pi^2$ 项给出
\begin{equation}
 \boxed{m_\pi^2=\xi m_A^2.}
 \label{eq:rxi-goldstone-mass-abelian-dp42}
\end{equation}
因此 $\pi$ 在固定规范中的质量依赖 $\xi$，不能把它解释成独立可观测粒子的质量；可观测散射矩阵中的规范参数依赖必须在完整规范恒等式中相消。

对外部有质量规范玻色子的物理四动量 \(k^\mu\)，若 \(E\gg m_Vc^2\)，纵向偏振满足
\begin{equation*}
 \epsilon_L^\mu(k)
 =\frac{k^\mu}{m_Vc}
 +\order\!\left(\frac{m_Vc^2}{E}\right).
\end{equation*}
Slavnov--Taylor 恒等式把 \(k_\mu\mathcal M^\mu\) 与对应 Goldstone 振幅联系起来，因此
\begin{equation}
 \mathcal M(\cdots V_L\cdots)
 =C\,\mathcal M(\cdots\pi\cdots)
 +\order\!\left(\frac{m_Vc^2}{E}\right),
 \label{eq:goldstone-equivalence-general-dp11}
\end{equation}
其中常用壳上树级约定给 \(C=1\)。控制量为 \(\epsilon_{\rm GE}=m_Vc^2/E\)。当 \(\epsilon_{\rm GE}\not\ll1\)（尤其接近产生阈值）时，高能展开失效，必须使用完整纵向偏振矢量与规范一致的原振幅；高阶修正还会改变有限的匹配因子 \(C\)。

整体破缺生成元通过守恒荷要求真实无质量 Goldstone 模；局域规范冗余经规范固定与 BRST 量子化后，把真空切向自由度与纵向规范偏振重组为有质量矢量表示。一般质量矩阵\eqnrefs{eq:general-gauge-boson-mass-matrix-dp11} 的零空间给出未破缺子代数，非零本征值给出规范玻色子质量。

对于一般多分量标量场，前面的圆形真空可以写成一个不依赖具体表示的代数结论。若势在生成元 \(T^a\) 下不变，则
\begin{equation}
\boxed{
\frac{\partial V}{\partial\phi_i}(T^a)_{ij}\phi_j=0.}
\label{eq:ssb-general-potential-invariance-dp68}
\end{equation}
在平稳真空 \(v\) 处再对场求导，定义 Hessian \(H_{ki}(v)=\partial_k\partial_iV|_v\)，得到
\begin{equation}
\boxed{
H_{ki}(v)(T^av)_i=0.}
\label{eq:ssb-general-hessian-zero-mode-dp68}
\end{equation}

\begin{proposition}[不变势在平稳点的 Hessian 零方向]
设实势函数二阶可微并在连续群作用下不变，$v$ 为平稳点，$T_Av$ 为群轨道的非零切向量。则有以下结论（原式 \eqnrefs{eq:ssb-general-hessian-zero-mode-dp68}）：
\begin{equation*}
\boxed{
H_{ki}(v)(T^av)_i=0.}
\end{equation*}
\end{proposition}
\begin{proof}
把一组实标量场坐标记为 $\phi_i$。连续整体对称性的无穷小作用写成

\begin{equation*}
 \delta\phi_i=\epsilon^a(T^a)_{ij}\phi_j,
\end{equation*}
其中 $\epsilon^a$ 是无穷小实参数，$T^a$ 是该实场表示中的生成元。势能在整个场空间上保持不变，因此对任意场构型都有
\begin{align*}
 0=\delta V
 &=\frac{\partial V}{\partial\phi_i}\,\delta\phi_i,\\
 &=\epsilon^a\frac{\partial V}{\partial\phi_i}(T^a)_{ij}\phi_j.
\end{align*}
由于各 $\epsilon^a$ 相互独立，每个生成元分别满足函数恒等式
\begin{equation*}
 \frac{\partial V}{\partial\phi_i}(T^a)_{ij}\phi_j=0.
\end{equation*}
对 $\phi_k$ 再求一次偏导，乘积法则给
\begin{equation*}
 0=
 \frac{\partial^2V}{\partial\phi_k\partial\phi_i}
 (T^a)_{ij}\phi_j
 +\frac{\partial V}{\partial\phi_i}(T^a)_{ik}.
\end{equation*}
令 $v_i$ 为平稳真空，即
\begin{equation*}
 \left.\frac{\partial V}{\partial\phi_i}\right|_{\phi=v}=0.
\end{equation*}
第二项在真空处消失，于是
\begin{align*}
 0
 &=\left.\frac{\partial^2V}{\partial\phi_k\partial\phi_i}
 \right|_{v}(T^av)_i,\\
 &=H_{ki}(v)(T^av)_i,
\end{align*}
其中
\begin{equation*}
 H_{ki}(v)\equiv
 \left.\frac{\partial^2V}{\partial\phi_k\partial\phi_i}\right|_v
\end{equation*}
是势能在真空处的 Hessian 矩阵。若生成元 $T^a$ 被真空自发破缺，则 $T^av\neq0$；上式说明非零向量 $T^av$ 属于 $H(v)$ 的零空间。对具有标准正定动能项的实标量场，二次拉氏量中的质量平方矩阵正比于该 Hessian 矩阵，所以这个真空轨道切向方向对应零质量模。Goldstone 零模因此来自“势能沿连续对称轨道保持常数”这一一般几何事实，而不依赖墨西哥帽势的具体形状。
\end{proof}

因此只有当 \(T^av\neq0\) 时，这个生成元才给出一个非平凡 Hessian 零方向。







\begin{derivation}[对称破缺后的径向与角向质量]

正文\eqnrefs{eq:ssb-complex-scalar-lag-dp11} 给出了只依赖 \(|\varphi_{\rm ssb}|\) 的对称势。先由驻点条件确定真空半径，再在所选真空附近展开到二次阶；径向与切向二次系数的差别最终给出正文\eqnrefs{eq:ssb-higgs-goldstone-masses-dp11} 的有质量径向模和无质量切向模。


令 \(r^2\equiv\varphi_{\rm ssb}^*\varphi_{\rm ssb}\)，则
\begin{equation*}
 V(r)=-\mu_{\rm ssb}^2r^2+\frac\lambda{\hbar c}r^4.
\end{equation*}
极值条件为
\begin{align*}
 \frac{\dd V}{\dd r}
 &=-2\mu_{\rm ssb}^2r+\frac{4\lambda}{\hbar c}r^3,\\
 &=2r\left(-\mu_{\rm ssb}^2+\frac{2\lambda}{\hbar c}r^2\right)=0.
\end{align*}
除去不稳定的 \(r=0\)，稳定极小值满足
\begin{equation*}
 r^2=\frac{\mu_{\rm ssb}^2\hbar c}{2\lambda}.
\end{equation*}
定义 \(r=v/\sqrt2\) 即得到
\begin{equation*}
 v^2=\frac{\mu_{\rm ssb}^2\hbar c}{\lambda}.
\end{equation*}

将真空附近参数化代入
\begin{equation*}
 \varphi_{\rm ssb}=\frac1{\sqrt2}(v+h+\ii\pi),
 \qquad
 \varphi_{\rm ssb}^*\varphi_{\rm ssb}=\frac12[(v+h)^2+\pi^2].
\end{equation*}
保留到二次阶：
\begin{align*}
 V
 ={}&-\frac{\mu_{\rm ssb}^2}{2}
 (v^2+2vh+h^2+\pi^2)\\
 &+\frac\lambda{4\hbar c}
 (v^2+2vh+h^2+\pi^2)^2.
\end{align*}
线性 \(h\) 项系数为
\begin{equation*}
 -\mu_{\rm ssb}^2v+\frac{\lambda v^3}{\hbar c}=0,
\end{equation*}
二次项为
\begin{equation*}
 V^{(2)}=\mu_{\rm ssb}^2h^2+0\times\pi^2.
\end{equation*}
与 \(\tfrac12(m^2c^2/\hbar^2)\phi^2\) 比较即得\eqnrefs{eq:ssb-higgs-goldstone-masses-dp11}。
\end{derivation}

\begin{derivation}[规范固定怎样消去 Goldstone 混合项]
直角坐标展开给出
\[
\Lag_{\rm mix}=g_{\rm G}v\,\mathscr C^\mu\partial_\mu\pi,
\qquad
F=\partial\!\cdot\!\mathscr C-\xi\frac{g_{\rm G}v}{\hbar c}\pi .
\]
规范固定项展开为
\[
\Lag_{\rm gf}
=-\frac{\hbar c}{2\xi}(\partial\!\cdot\!\mathscr C)^2
+g_{\rm G}v\,\pi\,\partial\!\cdot\!\mathscr C
-\frac{\xi g_{\rm G}^2v^2}{2\hbar c}\pi^2 .
\]
对中间项分部积分，
\[
g_{\rm G}v\,\pi\,\partial_\mu\mathscr C^\mu
=-g_{\rm G}v\,\mathscr C^\mu\partial_\mu\pi+\text{边界项},
\]
故它与 $\Lag_{\rm mix}$ 相消。剩余 $\pi^2$ 项与标准质量项比较：
\[
-\frac{\xi g_{\rm G}^2v^2}{2\hbar c}\pi^2
=-\frac12\left(\frac{m_\pi c}{\hbar}\right)^2\pi^2,
\qquad
\left(\frac{m_Ac}{\hbar}\right)^2=\frac{g_{\rm G}^2v^2}{\hbar c},
\]
于是
\[
\boxed{m_\pi^2=\xi m_A^2}.
\]
$\pi$ 的质量依赖规范参数 $\xi$，而物理规范玻色子质量 $m_A$ 不随 $\xi$ 改变。
\end{derivation}

\begin{derivation}[标量四次相互作用的一圈泡积分]
\eqnrefs{eq:dim-reg-measure-general-dp10} 中显式 \(\hbar\) 的消去可由变量替换直接看出。若从逆长度圈变量 \(\kappa^\mu\) 开始，定义物理动量

\begin{equation*}
 \ell^\mu=\hbar\kappa^\mu,
 \qquad
 \mu_{\rm DR}=\hbar\mu_\kappa,
\end{equation*}
则
\begin{equation*}
 \dd^d \kappa=\frac{\dd^d \ell}{\hbar^d},
 \qquad
 (\kappa^2-M^2)^{-2}
 =\frac{\hbar^4}{(\ell^2-m_R^2c^2)^2},
\end{equation*}
其中 \(M=m_Rc/\hbar\)。维数正则化尺度给
\begin{equation*}
 \mu_\kappa^{2\epsilon}
 =\frac{\mu_{\rm DR}^{2\epsilon}}{\hbar^{2\epsilon}}.
\end{equation*}
由于 \(d=4-2\epsilon\)，总的 \(\hbar\) 次数为
\begin{equation*}
 4-d-2\epsilon=0.
\end{equation*}
因此变量换元后的物理四动量形式为
\begin{equation*}
 \mu_{\rm DR}^{2\epsilon}
 \int\frac{\dd^d \ell}{(2\pi)^d}
 \frac{\ii}{(\ell^2-m_R^2c^2+\ii0^+)
 ((\ell+p)^2-m_R^2c^2+\ii0^+)},
\end{equation*}
其中显式 $\hbar$ 已由测度、分母和正则化尺度的幂次精确抵消。

用 Feynman 参数
\begin{equation*}
 \frac1{AB}=\int_0^1\frac{\dd x}{[xA+(1-x)B]^2}
\end{equation*}
并令 \(L=\ell+xp\)，分母化为
\begin{equation*}
 L^2-\Delta_s(x)+\ii0^+,
 \qquad
 \Delta_s(x)=m_R^2c^2-x(1-x)s.
\end{equation*}
于是
\begin{equation*}
 I(s)=\int_0^1\dd x\,
 \mu_{\rm DR}^{2\epsilon}
 \int\frac{\dd^d L}{(2\pi)^d}
 \frac{\ii}{(L^2-\Delta_s+\ii0^+)^2}.
\end{equation*}
做 Wick 转动 \(L^0=\ii L_E^0\)：
\begin{equation*}
 \dd^d L=\ii\dd^d L_E,
 \qquad
 L^2-\Delta+\ii0^+=-(L_E^2+\Delta-\ii0^+),
\end{equation*}
故
\begin{equation*}
 I(s)=-\int_0^1\dd x\,
 \mu_{\rm DR}^{2\epsilon}
 \int\frac{\dd^d L_E}{(2\pi)^d}
 \frac1{(L_E^2+\Delta_s-\ii0^+)^2}.
\end{equation*}
利用欧氏母积分
\begin{equation*}
 \int\frac{\dd^d L_E}{(2\pi)^d}
 \frac1{(L_E^2+\Delta)^2}
 =\frac1{(4\pi)^{d/2}}
 \Gamma\!\left(2-\frac d2\right)
 \Delta^{d/2-2},
\end{equation*}
代入 \(d=4-2\epsilon\)：
\begin{equation*}
 I(s)=-\int_0^1\dd x\,
 \frac1{(4\pi)^2}
 \Gamma(\epsilon)(4\pi)^\epsilon
 \left(\frac{\mu_{\rm DR}^2}{\Delta_s-\ii0^+}\right)^\epsilon.
\end{equation*}
逐项展开
\begin{align*}
 \Gamma(\epsilon)&=\frac1\epsilon-\gamma_E+\order(\epsilon),\\
 (4\pi)^\epsilon&=1+\epsilon\ln4\pi+\order(\epsilon^2),\\
 \left(\frac{\mu_{\rm DR}^2}{\Delta_s-\ii0^+}\right)^\epsilon
 &=1+\epsilon\ln\frac{\mu_{\rm DR}^2}{\Delta_s-\ii0^+}
 +\order(\epsilon^2).
\end{align*}
相乘到有限阶并执行参数积分，得到
\begin{equation*}
 I(s)=-\frac1{16\pi^2}
 \left[
 \frac1{\bar\epsilon}
 -\int_0^1\dd x\,
 \ln\frac{\Delta_s(x)-\ii0^+}{\mu_{\rm DR}^2}
 \right]+
 \order(\epsilon),
\end{equation*}
即\eqnrefs{eq:phi4-bubble-result-dp10}。当 \(\Delta_s(x)\) 穿过零时，\(-\ii0^+\) 决定对数的虚部，并与在壳中间态的 Cutkosky 不连续度一致。
\end{derivation}

\begin{derivation}[标量四次耦合的一圈演化]
记

\begin{equation*}
 a_1\equiv\frac{3}{32\pi^2},
 \qquad
 P_\epsilon\equiv\frac1{\bar\epsilon}.
\end{equation*}
一圈阶裸耦合为
\begin{equation*}
 g_{4,B}
 =\mu^{2\epsilon}
 \left[g_4(\mu)+a_1P_\epsilon g_4^2(\mu)
 +\order(g_4^3)\right].
\end{equation*}
按正文\eqnrefs{eq:phi4-beta-one-loop-dp10} 中 $\beta_{g_4}$ 的定义，$d=4-2\epsilon$ 时固定裸耦合的总尺度导数写成
\begin{equation*}
 \mu\frac{\dd g_4}{\dd\mu}\Big|_{g_{4,B}}
 =-2\epsilon g_4+\beta_{g_4}(g_4).
\end{equation*}
对裸表达式求导时，反项同样通过 \(g_4(\mu)\) 依赖于 \(\mu\)：
\begin{align*}
 0={}&2\epsilon
 \left(g_4+a_1P_\epsilon g_4^2\right)\\
 &+\left(1+2a_1P_\epsilon g_4\right)
 \left(-2\epsilon g_4+\beta_{g_4}\right)
 +\order(g_4^3).
\end{align*}
保留到 \(g_4^2\) 后，\(2\epsilon g_4\) 项相消，并得到
\begin{equation*}
 0=\beta_{g_4}
 -2a_1\epsilon P_\epsilon g_4^2
 +\order(g_4^3,\epsilon g_4^2).
\end{equation*}
因为
\begin{equation*}
 \epsilon P_\epsilon
 =\epsilon\left(\frac1\epsilon-\gamma_E+\ln4\pi\right)
 \longrightarrow1,
\end{equation*}
四维极限给出
\begin{equation*}
 \beta_{g_4}=2a_1g_4^2+
\order(g_4^3)
 =\frac{3g_4^2}{16\pi^2}+
\order(g_4^3).
\end{equation*}
这里关键的因子 2 来自 \(d=4-2\epsilon\) 下的 \(\mu^{2\epsilon}\)；若把反项对 \(g_4(\mu)\) 的隐含尺度依赖漏掉，就会错误地得到相反符号。
\end{derivation}

\begin{derivation}[固有时积分中的一圈紫外发散]

对正定欧几里得二次算符 $\Delta$，一圈实玻色 高斯行列式可用正文已经建立的固有时表示写成
\[
\frac12\Tr\ln\Delta
=-\frac12\mu_{\rm DR}^{2\epsilon}
\int_0^\infty\frac{\dd s}{s}\,
\Tr \ee^{-s\Delta},
\]
其中与场无关的加法常数不影响局域反项。把正文\eqnrefs{eq:heat-kernel-a4-flat-dp11} 代入，并只保留含 $s^2$ 的四维局域系数，得到
\begin{align*}
\left.\frac12\Tr\ln\Delta\right|_{a_4}
={}&-\frac12\mu_{\rm DR}^{2\epsilon}
\int\dd^d x\,
\tr a_4(x)\\
&\times\int_0^\infty\frac{\dd s}{s}
\frac{s^2}{(4\pi s)^{d/2}}.
\end{align*}
取 $d=4-2\epsilon$ 后，短固有时端的幂次变为
\[
\frac1s\frac{s^2}{s^{2-\epsilon}}
=s^{-1+\epsilon}.
\]
因此紫外极点只需从 $s\to0^+$ 端读取。任取固定分界尺度 $s_0$，
\begin{align*}
\mu_{\rm DR}^{2\epsilon}
\int_0^{s_0}\dd s\,s^{-1+\epsilon}
&=\mu_{\rm DR}^{2\epsilon}
\frac{s_0^\epsilon}{\epsilon}\\
&=\frac1\epsilon
+\ln(\mu_{\rm DR}^2s_0)+O(\epsilon).
\end{align*}
所以极点系数与任意分界 $s_0$ 无关，而长固有时部分只改变有限项。代入
\[
 a_4(x)
 =\frac12E^2
 +\frac1{12}\Omega_{\mu\nu}\Omega^{\mu\nu}
 +\frac16\nabla^2E
\]
并使用无边界体积分
\(
\int\dd^4 x\,\nabla^2E=0
\)，得到
\[
\left.\frac12\Tr\ln\Delta\right|_{\rm pole}
=-\frac1{2(4\pi)^2}\frac1\epsilon
\int\dd^4 x\,
\tr\left[
\frac12E^2+\frac1{12}\Omega_{\mu\nu}\Omega^{\mu\nu}
\right].
\]
最后使用正文已经固定的 $1/\bar\epsilon$ 记号整理 $1/\epsilon-\gamma_E+\ln4\pi$，即得到\eqnrefs{eq:heat-kernel-divergence-master-dp11}。
\end{derivation}

\begin{derivation}[Abel Higgs 模型中的矢量质量]
从
\[
\varphi_{\rm ssb}(x)
=\frac{v+h(x)}{\sqrt2}\ee^{\ii\pi(x)/v},
\qquad
D_\mu=\partial_\mu+\ii g_{\rm G}\mathscr C_\mu
\]
出发。局域 $U(1)$ 变换
\[
\varphi_{\rm ssb}\to \ee^{\ii\alpha(x)}\varphi_{\rm ssb},
\qquad
\mathscr C_\mu\to
\mathscr C_\mu-\frac1{g_{\rm G}}\partial_\mu\alpha
\]
要求极坐标变量变成
\[
h\to h,
\qquad
\pi\to\pi+v\alpha.
\]
因此组合
\[
\mathcal B_\mu
\equiv
\mathscr C_\mu+\frac1{g_{\rm G}v}\partial_\mu\pi
\]
满足 $\mathcal B_\mu\to\mathcal B_\mu$，它才是破缺相位中的规范不变矢量变量。

显式求协变导数：
\begin{align*}
D_\mu\varphi_{\rm ssb}
&=\frac{\ee^{\ii\pi/v}}{\sqrt2}
\left[
\partial_\mu h
+\ii(v+h)
\left(\frac{\partial_\mu\pi}{v}+g_{\rm G}\mathscr C_\mu\right)
\right]\\
&=\frac{\ee^{\ii\pi/v}}{\sqrt2}
\left[
\partial_\mu h
+\ii g_{\rm G}(v+h)\mathcal B_\mu
\right].
\end{align*}
取模平方时，实的径向项与纯虚的相位项交叉项严格相消：
\[
(D_\mu\varphi)^*(D^\mu\varphi)
=\frac12(\partial h)^2
+\frac{g_{\rm G}^2}{2}(v+h)^2\mathcal B_\mu\mathcal B^\mu.
\]
展开 $(v+h)^2=v^2+2vh+h^2$。与矢量二次型有关的纯二次项是
\[
\frac{g_{\rm G}^2v^2}{2}\mathcal B_\mu\mathcal B^\mu.
\]
在上述国际单位制归一化下，Proca 质量项写成
\[
\frac{\hbar c}{2}
\left(\frac{m_Ac}{\hbar}\right)^2
\mathcal B_\mu\mathcal B^\mu.
\]
两者是同一个二次算符，所以系数必须相等：
\[
g_{\rm G}^2v^2
=\hbar c\left(\frac{m_Ac}{\hbar}\right)^2,
\]
从而
\[
\boxed{m_Ac^2=g_{\rm G}\sqrt{\hbar c}\,v.}
\]
Goldstone 场 $\pi$ 通过 $\mathcal B_\mu$ 重组为有质量矢量场的纵向偏振。破缺前自由度计数为“复标量的两个实自由度 + 无质量矢量的两个横向偏振”共四个；破缺后为“径向标量 $h$ 一个 + 有质量矢量 $\mathcal B_\mu$ 三个偏振”仍共四个。这个计数是上述变量重组没有丢失物理自由度的独立检查。
\end{derivation}

